\documentclass[a4paper]{article}
\usepackage[left=10truemm,right=10truemm,top=25truemm,bottom=20truemm]{geometry}
\usepackage{mathtools}
\usepackage{amsmath}
\usepackage{amsfonts}
\usepackage{bm}
\usepackage{setspace}
\usepackage{wrapfig}
\usepackage{docmute}
\usepackage{amssymb}
\usepackage{hyperref}
\DeclareMathOperator\arctanh{arctanh}
\DeclareMathOperator\arccosh{arccosh}
\newcommand{\Jnorm}{J_{\rm norm}}
\newcommand{\Mnorm}{M_{\rm norm}}
\newcommand{\Jbound}{\frac{3\pi^2}{8}}

\title{Double Spacetime Theory:\\ Gravity as Torsion between Usual and Dual Spacetime}

\author{https://github.com/hypernumbernet}
\date{\today}

\begin{document}

\maketitle

\begin{abstract}
General Relativity has taught us to see gravity as the curvature of a single spacetime continuum. That teaching accounts for the motion of planets, the delay of light, and the ringing of black holes. It does not say why a continuum should curve, nor why the same phenomenon should appear as inertia under acceleration. The present work answers those questions by refusing the continuum as a primitive. Spacetime is not an autonomous stage. Each massive particle carries a pair of Minkowski frames---the usual, in which we live and measure, and a dual, time-reversed relative to the first, with the same interval. Gravity is the torsion of that pair: a mismatch that belongs to the particle itself.

From this picture two things follow at once. The geometry outside a spherical mass is recovered as a mismatch of the two frames, without taking curvature as a primitive. Gravity and inertia become the same phenomenon: free fall is the state in which the two frames stay aligned, and acceleration is the same misalignment when it is sourced by surrounding mass. If the ontology is sound, then the geometry of the world is not prior to matter. It is an attribute of particles, as spin and charge are attributes, and the apparent curvature of the heavens is the collective torsion of their dual frames.
\end{abstract}

\section{Introduction}
\label{sec:introduction}

Since 1915, general relativity has been the most successful description of gravitation we possess. The perihelion of Mercury, the delay of light, gravitational waves, and the imaging of black holes all speak for it. Yet a philosophical unease remains: why does matter curve spacetime, and why does curved spacetime dictate the motion of matter? The Einstein equations tell us how gravity works with exquisite precision. They do not say why a continuum should curve at all. The unease sharpens against Mach's principle. The inertia of a body ought to be determined by the totality of matter in the universe, yet general relativity treats spacetime as an independent continuum that can exist even in complete emptiness.

A century of attempts to close this gap has left the continuum itself untouched. Spacetime is taken to be a smooth manifold a priori, and gravity is the metric's deviation from flatness.

The present work refuses that primitive. There is no single shared spacetime. Each massive particle carries its own pair of Minkowski frames. One is the usual spacetime, in which we live and measure. The other is a dual spacetime: time-reversed relative to the first, with the same interval. Gravity and inertia are not curvature of a background. They are the mismatch of these two frames when one particle is asked to stay parallel to another.

The shift is ontological. In the standard picture, spacetime tells matter how to move and matter tells spacetime how to curve. Here there is no autonomous spacetime at all. The degrees of freedom we call spacetime are internal attributes of particles, as spin and charge are attributes. The apparent curvature of the macroscopic world is the collective mismatch of neighboring particles.

Two consequences follow at once, and they can be stated before any formula is written. First, the geometry outside a spherical mass is recovered as a mismatch of the two frames, without taking curvature as a primitive. Second, gravity and inertia are the same phenomenon. Free fall is the state in which the two frames stay aligned. Acceleration breaks that alignment, and the resulting mismatch is the inertial force---the same torsion that we call gravity when it is sourced by surrounding mass.

This is not a modification of general relativity. It is a proposed completion: the continuum remains a useful macroscopic fiction, while spacetime degrees of freedom are treated as paired and particle-local.

The paper proceeds in that order. Sections~\ref{sec:math} and~\ref{sec:kinematics} introduce the two sectors and the transformations that move them. Section~\ref{sec:gravity} defines gravity as their mismatch and recovers the exterior of a spherical mass. Section~\ref{sec:unified} reads inertia from the same mismatch, so that the equivalence principle is no longer a postulate. With that picture in hand, later sections take the geometry to strong fields, to nuclear and atomic scales, to electromagnetism, and to galaxies. Those readings are built on the picture; they are not its foundation. Section~\ref{sec:conclusions} returns to the ontology and to the observations that would test it.

\section{Mathematical Foundation: The 16-Dimensional Biquaternion Algebra}
\label{sec:math}

The entire double spacetime theory is constructed within the 16-real-dimensional algebra of biquaternions (also known as double quaternions), which is canonically isomorphic to the Clifford geometric algebra $\mathrm{Cl}(3,1)$ over Minkowski spacetime with signature $(-,+,+,+)$. This algebra naturally accommodates two complete, mutually commuting copies of Minkowski spacetime — the usual and the dual — within a single algebraic structure intrinsically attached to every massive particle.

\subsection{The Biquaternion Algebra}

We begin with two independent copies of Hamilton's quaternion algebra:
\begin{itemize}
  \item Primary quaternions generated by $i,j,k$ obeying $i^2=j^2=k^2=-1$, \ $ij=k$, \ $ji=-k$ (cyclic),
  \item Secondary quaternions generated by $I,J,K$ obeying identical rules $I^2=J^2=K^2=-1$, \ $IJ=K$, \ $JI=-K$ (cyclic).
\end{itemize}
The full biquaternion algebra is their tensor product $\mathbb{H} \otimes_{\mathbb{R}} \mathbb{H}$, in which the two sets strictly commute:
\[
iI = Ii,\; iJ = Ji,\; iK = Ki,\; jI = Ij,\; \dots
\]
The resulting 16 linearly independent real basis elements are
\[
1,\; i,\; j,\; k,\; I,\; J,\; K,\; iI,\; iJ,\; iK,\; jI,\; jJ,\; jK,\; kI,\; kJ,\; kK.
\]

\subsection{Explicit Isomorphism with Cl(3,1)}

The algebra is isomorphic to $\mathrm{Cl}(3,1)$ via the following identification of grade-1 (vector) generators:
\begin{align*}
  e_0 &= j,                    & &\text{(time-like basis vector)}\\[4pt]
  e_1 &= kI, \quad e_2 = kJ, \quad e_3 = kK & &\text{(space-like basis vectors)},
\end{align*}
which immediately satisfy the defining relations
\[
e_0^2 = -1, \quad e_1^2 = e_2^2 = e_3^2 = +1, \quad \{e_\mu,e_\nu\} = 0 \;\; (\mu \neq \nu).
\]
The higher-grade elements then follow standard Clifford multiplication. In particular:
\begin{alignat*}{3}
  &\text{grade-2 (bivectors):}  &\quad& e_0e_1 = iI, \;\; e_0e_2 = iJ, \;\; e_0e_3 = iK, \\
  &&& e_3e_2 = I,   \;\; e_1e_3 = J,   \;\; e_2e_1 = K, \\[6pt]
  &\text{grade-3 (trivectors):} && e_1e_2e_3 = k, \;\; e_0e_3e_2 = jI, \;\; e_0e_1e_3 = jJ, \;\; e_0e_2e_1 = jK, \\[6pt]
  &\text{pseudoscalar:}        && e_0e_1e_2e_3 = i.
\end{alignat*}
The complementary ``dual'' basis (time-reversed under right-multiplication by $i$) is generated by
\[
\tilde{e}_0 = k = e_1e_2e_3, \quad \tilde{e}_1 = jI = e_0e_3e_2, \quad \tilde{e}_2 = jJ = e_0e_1e_3, \quad \tilde{e}_3 = jK = e_0e_2e_1,
\]
yielding the dual spacetime representation introduced below.

This explicit isomorphism makes manifest that the 16-dimensional biquaternion algebra is not an ad-hoc construction but the geometrically natural Clifford algebra of 4-dimensional Minkowski spacetime, now enriched with an intrinsic duality that distinguishes usual and dual sectors while preserving the full Lorentzian structure. 

The apparent irregularity in the Clifford algebraic forms — such as the asymmetric placement of $j$ as the sole time-like vector generator and the cross-term-dominated spatial basis $kI$, $kJ$, $kK$ — is not accidental: it exists precisely to preserve the full rotational symmetry of three-dimensional physical space within a structure that simultaneously encodes two time-reversed copies of Minkowski spacetime. This subtle asymmetry is what allows the dual map $X \mapsto Xi$ to reverse the arrow of time without breaking spatial isotropy, thereby providing the algebraic foundation for both the equivalence principle and the emergence of gravitational torsion as a relative misalignment between particle-intrinsic dual frames.

\subsection{Dual Spacetime Structure Intrinsic to Particles}

Each particle is postulated to carry an intrinsic pair of compactified Minkowski spacetimes encoded within the biquaternion algebra:
\begin{itemize}
  \item \textbf{Usual Spacetime:} Represented by vectors of the form
  \[
  X = ct \, j + x \, kI + y \, kJ + z \, kK,
  \]
  with Minkowski norm
  \[
  X^2 = -(ct)^2 + x^2 + y^2 + z^2.
  \]

  \item \textbf{Dual Spacetime:} Represented by vectors of the form
  \[
  X' = ct' \, k + x' \, jI + y' \, jJ + z' \, jK,
  \]
  with identical Minkowski norm
  \[
  X'^2 = -(ct')^2 + x'^2 + y'^2 + z'^2.
  \]
\end{itemize}
The two spacetimes are related by the dual map
\[
X' = X i,
\]
which reverses the arrow of time:
\[
(ct \, j) i = ct \, (j i) = -ct \, k,
\]
while preserving the Minkowski norm. Since the pseudoscalar $i = e_0 e_1 e_2 e_3$ of $\mathrm{Cl}(3,1)$ anticommutes with every grade-1 vector ($i X = -X i$ for $X \in \mathrm{Cl}(3,1)$) and satisfies $i^2 = -1$, the norm transforms as
\[
X'^2 = (X i)(X i) = X (i X) i = -X (X i) i = -X^2\, i^2 = -X^2 \cdot (-1) = X^2.
\]
This dual structure is intrinsic to every particle, with the usual spacetime governing its standard kinematics and the dual spacetime encoding complementary degrees of freedom that become dynamically relevant in the presence of gravity and inertia. The existence of these two spacetimes allows for a richer geometric interpretation of Lorentz transformations, parallel transport, and ultimately the nature of gravitation itself.

\section{Kinematic Structure: Rotors, Boosts, and Parallel Transport}
\label{sec:kinematics}

In this section we develop the kinematic framework of the theory. Lorentz boosts of the usual sector and elliptic rotations of the dual sector are generated by the six bivectors $\{i\Gamma_a,\Gamma_a\}$. Same-axis pairs commute; off-axis pairs need not. Boosts arise from bivectors that square to $+1$, yielding the $\cosh$/$\sinh$ expansion of the rotor.

\subsection{Rotors in the Usual Spacetime}

To illustrate the kinematic structure, consider a Lorentz boost along the $x$-direction with rapidity $\phi$. The generating bivector is $iI$, which satisfies $(iI)^2 = i^2 I^2 = (-1)(-1) = +1$. In general, the usual-sector boost generators $iI$, $iJ$, $iK$ square to $+1$, while the dual-sector rotation generators $I$, $J$, $K$ square to $-1$. The same $I$, $J$, $K$ are the spatial bivectors of the usual Minkowski frame ($e_3e_2=I$, $e_1e_3=J$, $e_2e_1=K$); they do not furnish a second, independent copy of $\mathfrak{so}(3)$.

The corresponding boost rotor is thus
\[
R_\text{usual} = \exp\!\left( \frac{\phi}{2} iI \right).
\]
Given the positive square of the generator, the exponential expands naturally into hyperbolic functions:
\[
R_\text{usual} = \cosh \frac{\phi}{2} + iI \sinh \frac{\phi}{2}.
\]
Applying this rotor via the sandwich product to a spacetime vector $X = ct \, j + x \, kI + \cdots$ yields the canonical Lorentz transformation:
\begin{align*}
ct' \, j + x' \, kI &= R_\text{usual} \, X \, R_\text{usual}^\dagger \\
&= \left[ \cosh\frac{\phi}{2} + iI \sinh\frac{\phi}{2} \right] (ct \, j + x \, kI) \left[ \cosh\frac{\phi}{2} - iI \sinh\frac{\phi}{2} \right] \\
&= (\gamma ct + \gamma v x) \, j + (\gamma x + \gamma v ct) \, kI,
\end{align*}
where $\gamma = \cosh \phi$ and $v/c = \tanh \phi$. Spatial rotations of the usual Minkowski vectors are generated by the same $\Gamma_a\in\{I,J,K\}$ that generate dual-sector elliptic rotors; they are not an extra Lorentz copy.

\subsection{The Complete Rotor}

The most general six-parameter exponential on this basis is
\[
R_\text{total} = \exp\left( \frac{\phi_1}{2} iI + \frac{\phi_2}{2} iJ + \frac{\phi_3}{2} iK + \frac{\theta_1}{2} I + \frac{\theta_2}{2} J + \frac{\theta_3}{2} K \right),
\]
where the parameters $\phi_a$ govern usual-sector boosts and $\theta_a$ dual-sector rotations. Write $\Gamma_1=I$, $\Gamma_2=J$, $\Gamma_3=K$ for the spatial bivectors of Section~\ref{sec:math}. Factorisation $R_{\mathrm{total}}=R_{\mathrm{usual}}R_{\mathrm{dual}}$ is not automatic: it requires the two summands to commute. About a single axis they do: $[i\Gamma_a,\Gamma_a]=0$, so
\[
\exp\Bigl(\tfrac{\phi}{2}i\Gamma_a+\tfrac{\theta}{2}\Gamma_a\Bigr)
=\exp\bigl(\tfrac{\phi}{2}i\Gamma_a\bigr)\,\exp\bigl(\tfrac{\theta}{2}\Gamma_a\bigr).
\]
About different axes they need not. In particular $[iI,J]\neq 0$, and the six-parameter exponential does not factor into a product of one-axis rotors without that extra assumption. In what follows we work on a single axis whenever we write $R_{\mathrm{total}}=R_{\mathrm{usual}}R_{\mathrm{dual}}$.

To make the geometric meaning more transparent on a single axis, we now introduce the bivector form of the rotor, which separates magnitudes from directions.

Define the unit boost bivector in the usual spacetime
\[
\hat{p} = p_1 \, iI + p_2 \, iJ + p_3 \, iK, \qquad \hat{p}^2 = +1,
\]
where $\hat{p}$ is normalized ($|\hat{p}| = 1$) and $\phi = \sqrt{\phi_1^2 + \phi_2^2 + \phi_3^2}$ is the total rapidity (with direction encoded in the components $p_a = \phi_a / \phi$).

Similarly, define the unit rotation bivector in the dual spacetime
\[
\hat{q} = q_1 \, I + q_2 \, J + q_3 \, K, \qquad \hat{q}^2 = -1,
\]
where $\hat{q}$ is normalized ($|\hat{q}| = 1$) and $\theta = \sqrt{\theta_1^2 + \theta_2^2 + \theta_3^2}$ is the total dual rotation angle (with direction encoded in $q_a = \theta_a / \theta$).

The complete rotor then takes the bivector-exponential form
\[
R_\text{total} = \exp\!\left( \frac{\phi}{2} \hat{p} + \frac{\theta}{2} \hat{q} \right).
\]
When $\hat{p}$ and $\hat{q}$ are aligned on a single axis the exponential factorises, $R_{\mathrm{total}}=R_{\mathrm{usual}}R_{\mathrm{dual}}$, with
\[
R_\text{usual} = \exp\!\left( \frac{\phi}{2} \hat{p} \right)
= \cosh\frac{\phi}{2} + \hat{p} \sinh\frac{\phi}{2},
\]
\[
R_\text{dual} = \exp\!\left( \frac{\theta}{2} \hat{q} \right)
= \cos\frac{\theta}{2} + \hat{q} \sin\frac{\theta}{2}.
\]
The six real parameters remain the rapidity $\phi$ and unit boost direction $\hat{p}$ in the usual sector, and the rotation angle $\theta$ and unit axis $\hat{q}$ in the dual sector. The factorisation is used below only when the generators commute.

\subsection{The Sandwich Product}

The versor formalism provides a compact means to apply these transformations: for any biquaternion multivector $X$ representing a spacetime event, the transformed vector is
\[
\tilde{X} = R_\text{total} \, X \, R_\text{total}^\dagger.
\]
This single operation preserves the Minkowski norm ($X^2 = \tilde{X}^2$) and acts coherently on both usual and dual components, seamlessly integrating boosts and rotations into one algebraic expression. No matrices or coordinate-dependent machinery are required; the geometry unfolds intrinsically within the algebra.

\section{Gravity as Torsional Mismatch between Dual Spacetime}
\label{sec:gravity}

In this central section, we establish that gravity is neither spacetime curvature nor an external force, but the torsional mismatch between the usual and dual spacetimes intrinsically carried by every massive particle. Christoffel symbols, Riemann tensors, and a continuous manifold are not primitive objects. Torsion is the relative rotation angle between particle-local dual rotors. The parameter-space scalar $J$ is distinct from the Weitzenb\"ock scalar $T$.

\subsection{Definition of the Torsional Mismatch Rotor}

The torsional mismatch rotor is defined as
\[
\Omega = R_\text{usual}^\dagger R_\text{dual}.
\]

The dagger operation is the standard conjugate-inverse of the rotor in the biquaternion algebra:
\[
R_\text{usual}^\dagger = \exp\!\left( -\frac{\phi}{2} \hat{p} \right)
= \cosh\frac{\phi}{2} - \hat{p} \sinh\frac{\phi}{2}.
\]
Although formally a simple conjugation, this dagger carries special significance in the double spacetime framework: the negative sign in front of the sinh term precisely encodes the time-reversal inherent in the dual map. Since the pseudoscalar $i$ is grade-4 and therefore commutes with bivectors, right-multiplication by $i$ acts on the boost generator as
\[
\hat{p}\, i = - (p_1 I + p_2 J + p_3 K) \equiv -\hat{q}',
\qquad (\hat{p}\, i)^2 = -1,
\]
where $\hat{q}' \equiv p_1 I + p_2 J + p_3 K$ is a unit rotation generator in the dual sector. Equivalently, $\hat{p} = \hat{q}'\, i$, and the inverse usual rotor admits the exact rewriting
\[
R_\text{usual}^\dagger = \cosh\frac{\phi}{2} - \hat{q}'\, i\, \sinh\frac{\phi}{2}.
\]
Using the standard identities $\cosh(\phi/2) = \cos(i\phi/2)$ and $\sinh(\phi/2) = \sin(i\phi/2)/i$, this admits the analytically continued ($\phi \to i\phi$) trigonometric form
\[
R_\text{usual}^\dagger = \cos\!\left(\frac{i\phi}{2}\right) - \hat{q}'\, \sin\!\left(\frac{i\phi}{2}\right),
\]
revealing that the dagger operation effectively bridges the hyperbolic geometry of the usual spacetime to the elliptic geometry of the dual spacetime---not as an equality of $\cosh$ and $\cos$ for real arguments, but as the unique analytic continuation of the boost rotor across the duality $X \mapsto X i$.

The relative rotor $\Omega = R_\text{usual}^\dagger R_\text{dual}$ realizes this bridge operationally: the dagger temporarily ``deconstructs'' the hyperbolic boost structure of the usual spacetime (inducing the dual map through its sign reversal), while subsequent multiplication by $R_\text{dual}$ ``reconstructs'' the combined object, making the relative misalignment manifest. Through this ``deconstruct-then-reconstruct'' process, the torsional mismatch between the usual and dual spacetimes is extracted algebraically, and this mismatch constitutes the essence of gravity in the double spacetime framework.

Free fall is the state in which this mismatch vanishes. Unitarity of $R_{\mathrm{usual}}$ gives a sharp criterion: $\Omega=1$ if and only if the two rotors coincide, $R_{\mathrm{dual}}=R_{\mathrm{usual}}$. Anti-synchronization $R_{\mathrm{dual}}=R_{\mathrm{usual}}^\dagger$ is a different condition; on a single axis it forces both bivector coefficients to vanish, so a nontrivial pure boost does not satisfy it. We read mass-energy as the physical cost of deviations from $\Omega=1$. The rotors themselves are units in the even subalgebra. The logarithm $\Omega_{\mathrm{biv}}=\log\Omega$ takes values in $\mathfrak{so}(3,1)$, and the six generators $\{i\Gamma_a,\Gamma_a\}$ span one Lorentz copy, not the twelve-dimensional sum $\mathfrak{so}(3,1)\oplus\mathfrak{so}(3,1)$.

The associated torsion bivector is
\[
\Omega_\text{biv} = \log \Omega \in \mathfrak{so}(3,1),
\]
the principal branch on a simply connected neighbourhood of the identity. Throughout this paper the regime of small mismatch is $\Omega \approx 1$, where $\Omega_\text{biv}$ is uniquely determined; topologically distinct branches lying outside this neighbourhood are not used in the dynamics presented here.

\subsection{Lorentz-Invariant Scalar from Dual Torsion}

The Killing form on $\mathfrak{so}(3,1)$ (vector representation) is
\[
B(X,Y) = 4\,\mathrm{Tr}(X Y).
\]
Denote the three grade-2 spatial bivectors of \S\ref{sec:kinematics} collectively by $\Gamma_a$ ($a=1,2,3$) with $\Gamma_1 = I$, $\Gamma_2 = J$, $\Gamma_3 = K$ (the Clifford identification $e_3e_2 = I$, $e_1e_3 = J$, $e_2e_1 = K$ of \S\ref{sec:math}). The usual-spacetime boost generators are $i\Gamma_a$, while $\Gamma_a$ generate rotations in the dual sector. Evaluating the Killing form gives $B(i\Gamma_a, i\Gamma_a) = +8$ and $B(\Gamma_a, \Gamma_a) = -8$, reproducing the Lorentzian signature perfectly. Same-axis generators commute, $[i\Gamma_a,\Gamma_a]=0$; off-axis commutators need not vanish.

Write the half-angle expansion $\Omega_\text{biv}=\sum_a\bigl(\frac{\alpha_a}{2}i\Gamma_a+\frac{\beta_a}{2}\Gamma_a\bigr)$. The six generators span one copy of $\mathfrak{so}(3,1)$; together with translations the ambient algebra is the Poincar\'e candidate $\mathfrak{iso}(3,1)\cong\mathfrak{so}(3,1)\ltimes\mathbb{R}^{3,1}$. Then
\[
B(\Omega_\text{biv},\Omega_\text{biv})=2\sum_{a=1}^3(\alpha_a^2-\beta_a^2).
\]
The scalar that measures the mismatch throughout this paper is
\begin{equation}
\label{eq:J-def}
J=\frac12\sum_{a=1}^3(\alpha_a^2-\beta_a^2),
\end{equation}
four times $\frac1{16}B(\Omega_\text{biv},\Omega_\text{biv})$ under this expansion. We also write $\Jnorm=\frac8{3\pi^2}J$, and $M=\frac12\sum_a(\alpha_a^2+\beta_a^2)$ for the unsigned torsional mass, with the same normalisation $\Mnorm=\frac8{3\pi^2}M$.

The sign of $J$ is the distinction between ordinary attractive gravity and dual-dominated repulsion. When usual-spacetime boosts dominate, $J>0$ and the force is attractive. When dual rotations dominate, $J<0$ and the force is repulsive. Swapping the two sectors, $\alpha\leftrightarrow\beta$, sends $J\mapsto -J$ and leaves $M$ invariant.

Three vanishing conditions must not be conflated. If $M=0$, then $J=0$: every rapidity vanishes, and this is vacuum. The converse is false. Equal usual and dual rapidities give $J=0$ with $M>0$---balanced massive states---and on a single axis those states still have $\Omega\neq 1$. Thus $\Omega=1$, $J=0$, and $M=0$ are three different physical situations. Gravity as free fall is $\Omega=1$, not the vanishing of $J$.

The algebra does not allow unbounded mismatch. On admissible configurations---non-negative rapidities with $\alpha_a+\beta_a\le\pi/2$ on each axis---one has $|J|\le\Jbound$ and $|\Jnorm|\le 1$, and likewise $0\le\Mnorm\le 1$. The image of $\Jnorm$ on this cone is the whole interval $[-1,1]$. Equality $|\Jnorm|=1$ holds only for a purely hyperbolic configuration ($\alpha_a=\pi/2$, $\beta_a=0$ on every axis) or a purely elliptic one ($\alpha_a=0$, $\beta_a=\pi/2$). A principal-branch condition $|\alpha_a+\beta_a|\le\pi/2$ without the sign restriction is not enough: the bound fails.

\subsection{\texorpdfstring{Schwarzschild Chart and the Distinction among $J$, $J_{\mathrm{field}}$, and $T$}{Schwarzschild Chart and J, J field, and T}}
\label{subsec:schwarzschild-chart}

On the exterior Schwarzschild domain $r>r_s>0$ the radial pure-boost rapidity is
\[
\varphi(r) = -\tfrac12\log\bigl(1-r_s/r\bigr).
\]
The associated sandwich scales reproduce the diagonal tetrad legs:
\[
\mathrm{e}^{-\varphi}=\sqrt{A},\qquad \mathrm{e}^{\varphi}=A^{-1/2},\qquad A(r)=1-r_s/r.
\]
The coframe $\bigl(\sqrt{A},\,A^{-1/2},\,r,\,r\sin\theta\bigr)$ induces the Schwarzschild metric $\mathrm{diag}(-A,\,A^{-1},\,r^2,\,r^2\sin^2\theta)$. More generally, a rotor field $R_{\mathrm{total}}(x)$ acts on a fixed reference tetrad $\{\mathring{e}^a\}$ by the sandwich product
\[
e^a{}_\mu(x) \;=\; \langle R_{\mathrm{total}}(x)\, \mathring{e}^a\, R_{\mathrm{total}}^\dagger(x) \rangle_{1,\mu},
\]
yielding the metric $g_{\mu\nu}=\eta_{ab}\,e^a{}_\mu e^b{}_\nu$. For a coframe $e^a{}_\mu$ the Weitzenb\"ock torsion is $T^a{}_{\mu\nu}=\partial_\mu e^a{}_\nu-\partial_\nu e^a{}_\mu$, and the teleparallel scalar is the quadratic combination
\[
T=\tfrac14 T_{\rho\mu\nu}T^{\rho\mu\nu}+\tfrac12 T_{\rho\mu\nu}T^{\nu\mu\rho}-T_\rho T^\rho.
\]
On the Schwarzschild coframe this evaluates to
\begin{equation}
\label{eq:schwarzschild-T}
T = \frac{2}{r^2}\frac{(1-\sqrt{A})^2}{\sqrt{A}} = -\frac{1}{r^2}\,\partial_r\bigl(r^2 B^r\bigr),\qquad B^r=\frac{4(1-\sqrt{A})}{r}.
\end{equation}
Three quadratic quantities appear on this chart and must not be conflated:
\[
J=\tfrac12\varphi^2,\qquad J_{\mathrm{field}}=\tfrac12(\varphi')^2,\qquad T\text{ as in }\eqref{eq:schwarzschild-T}.
\]
The first is the finite-angle mismatch of the two rotors; the second is a field seed built from $\partial_r\varphi$; the third is the teleparallel Lagrangian density. They are not the same object. A constant nonzero boost already has $J\neq 0$ with vanishing $T$ on flat space, so $J$ cannot be a multiple of $T$ in general. A dictionary for an arbitrary tetrad remains open; what follows is the dictionary on every static radial-boost chart.

Substituting $\sqrt{A}=\mathrm{e}^{-\varphi}$ into~\eqref{eq:schwarzschild-T} yields
\[
\frac{(1-\sqrt{A})^2}{\sqrt{A}}=\mathrm{e}^{\varphi}-2+\mathrm{e}^{-\varphi}=2(\cosh\varphi-1).
\]
Hence, for the time-leg $A>0$ of any static radial-boost coframe, with rapidity $\varphi=-\tfrac12\log A$ and $J=\tfrac12\varphi^2$,
\begin{equation}
\label{eq:jt-dictionary}
T=\frac{4}{r^2}(\cosh\varphi-1)=\frac{8}{r^2}\sinh^2\frac{\varphi}{2},
\qquad
r^2 T=4\bigl(\cosh\sqrt{2J}-1\bigr)=4J+\tfrac23 J^2+\tfrac{2}{45}J^3+\cdots.
\end{equation}
In the weak field, $\varphi\to 0$, one has $r^2 T/(4J)\to 1$: the two densities agree at leading order, and nowhere else. The map $J\mapsto r^2 T$ is strictly increasing on the whole admissible cone $|J|\le\Jbound$, so it inverts:
\[
J=\tfrac12\Bigl(\arccosh\bigl(1+\tfrac14 r^2 T\bigr)\Bigr)^{2}
\quad(J\ge 0),\qquad
J=-\tfrac12\Bigl(\arccos\bigl(1+\tfrac14 r^2 T\bigr)\Bigr)^{2}
\quad(J<0).
\]
On this gauge the teleparallel density at a given radius determines the algebraic scalar, even though the two are never equal outside the flat limit. Moreover $4J/r^2\le T$, with equality only in the flat case $A=1$; on the exterior Schwarzschild chart one has in addition $T\le 4J_{\mathrm{field}}$.

The analytic continuation $E(J)=\cosh\sqrt{2J}-1$ is entire. For $J<0$ it becomes $E(J)=\cos\sqrt{-2J}-1$. On the admissible cone the argument never reaches $\pi$, so $E(J)<0$ whenever $J<0$, and $\operatorname{sign} T=\operatorname{sign} J$. A real (hyperbolic) radial boost cannot produce $T<0$: on a static radial-boost gauge with $A>0$ one has $T\ge 0$, vanishing only for $A=1$. Repulsive layers live in the elliptic dual sector. The cosine branch is not a formal continuation; it is the geometry in which gravity reverses.

Among all static radial-boost gauges, the vacuum condition $\frac{d}{dr}(rA)=1$ selects a unique member: $A(r)=1-r_s/r$, with $r_s=r(1-A(r))$. That is the exterior Schwarzschild factor. On this chart, $\frac{r_s}{2r}\le\varphi(r)\le\frac{r_s}{2rA(r)}$, so $r\varphi\to r_s/2$ at infinity and the finite-angle rapidity tracks the Newtonian potential. The two densities share the same leading falloff, $r^4 T\to r_s^2/2$ and $4r^2 J\to r_s^2/2$. As $r\to r_s^+$ one has $A\to 0^+$, $\varphi\to\infty$, and $T\to\infty$: the classical chart is unbounded, a pathology the admissible cone will cut off in Section~\ref{subsec:quasi-horizon}.

\subsection{Proposed Action}

We propose the action
\[
S = \frac{c^4}{16\pi G} \int J \, d^4x
\]
as a Lagrangian density in the mismatch scalar $J$. The integral is algebraic and background-independent. It is not the teleparallel integral of $T$. For $T$ itself the standard identity
\[
\int T \, e \, d^4x = -\int R \, \sqrt{-g} \, d^4x + \int \partial_\mu(\sqrt{-g}\, V^\mu)\, d^4x
\]
holds after tetrad variation with vanishing boundary variations, so the Einstein equations belong to $T$, not to $J$. Gravity in this theory remains the mismatch angle between particle-intrinsic dual frames. On a static radial-boost coframe the teleparallel density dominates, $e\cdot 4J/r^2\le e\cdot T$, with equality only when $A=1$. The two agree in the weak field. Recovering the Einstein equations from $J$ for a general tetrad requires a dictionary that is not yet written.

\subsection{Weak-Field and Newtonian Limit}

In the linearized regime,
\[
R_\text{usual} \approx 1 + \frac{1}{2} \sum \alpha_a \, i\Gamma_a, \quad R_\text{dual} \approx 1,
\]
yielding $\Omega_\text{biv} \approx \sum \alpha_a \, i\Gamma_a$ and $J \approx \frac{1}{2} \sum \alpha_a^2$. If the boost parameters track the Newtonian potential as $\alpha_a \propto \partial_a (\Phi/c^2)$, Poisson's equation
\[
\nabla^2 \Phi = 4\pi G \rho
\]
is recovered with the standard coefficient. This matching identifies $G$ with the linearised chart; it does not derive $G$ from the algebra.

\subsection{Rotor Action in the Dual Spacetime}
\label{subsec:dual-rotors}

The sole algebraic degrees of freedom considered here are the rotor angles within $\mathbb{H} \otimes_{\mathbb{R}} \mathbb{H} \cong \mathrm{Cl}(3,1)$ carried by every massive particle. Gravity, inertia, and quantum states are read as deviations and coherences among these angles.

The dual spacetime, related to the usual spacetime by the dual map $X \mapsto Xi$, inherits a time-reversed orientation. Consequently, an angle that generates a pure spatial rotation in the dual sector must be appropriately dual-transformed when incorporated into the unified rotor description.

Consider a dual-sector angle $\theta$. In the pure dual basis, it produces a rotation generated by the elliptic versor $\hat{q}$ ($\hat{q}^2 = -1$):
\[
\exp\left( \frac{\theta}{2} \hat{q} \right).
\]
However, since the dual map reverses the arrow of time, the physical effect of this angle on the dual spacetime is equivalent to a boost when viewed through the dual transformation.

To reflect this correctly in the torsional mismatch calculation, we apply the dual conversion to the angle itself: $\theta \mapsto \theta i$. This transforms the dual contribution into a hyperbolic boost term:
\[
\exp\left( \frac{\theta i}{2} \hat{q} \right) = \exp\left( \frac{\theta}{2} \hat{q}i \right) = \exp\left( \frac{\theta}{2} (q_1 iI + q_2 iJ + q_3 iK) \right).
\]
Thus, the dual rotor is defined as
\[
R_\text{dual} = \exp\left( \frac{\theta}{2} \hat{q} \right),
\]
while its action on dual spacetime vectors, after the angle dual conversion $\theta \mapsto \theta i$, manifests as a boost generated by $i\Gamma_a$ ($ (\hat{q}i)^2 = +1 $).

In free fall, $\Omega=1$ requires $R_{\mathrm{dual}}=R_{\mathrm{usual}}$, not $J=0$. Mass-energy is the cost of deviations from that condition, and $J>0$ is attractive gravity.

This angle-centric treatment, with explicit dual conversion $\theta \mapsto \theta i$, preserves the original form of the total rotor
\[
R_\text{total} = \exp\left( \frac{\phi}{2} \hat{p} + \frac{\theta}{2} \hat{q} \right)
\]
while correctly incorporating the time-reversal duality: the dual rotor $R_\text{dual}$ generates rotations in its native basis, but its physical effect on the dual spacetime—through the dual-transformed angles—is that of a boost, naturally opposing the usual-sector boosts in vacuum.

\subsection{Physical Interpretation and Gravitational Engineering}

Gravity is the energetic cost of dual-rotor misalignment. Because $J$ depends solely on the relative rotor $\Omega$, an external influence that coherently modulates the dual components $\theta$ could in principle alter $J$, including the possibilities $J\to 0$ (shielding) and $J<0$ (repulsion). Whether such a coupling exists is a conjecture, taken up in Section~\ref{sec:gravitycontrol}. The continuum remains a macroscopic description; spacetime degrees of freedom in this theory are particle-local and paired.

\section{Unified Description of Inertia and Gravity}
\label{sec:unified}

In this section, we reveal the deepest consequence of the dual spacetime formalism: inertia and gravity are not merely equivalent — they are identical phenomena arising from the torsional mismatch between the usual and dual spacetimes. The distinction between inertial forces in accelerated frames and gravitational forces in curved spacetime dissolves completely. Both are the same dual rotor misalignment, with the equivalence principle emerging as a direct algebraic necessity rather than a postulate.

\subsection{Inertial Resistance as Dual Rotor Rigidity}

Rest mass $m > 0$ is interpreted as a coupling between the usual and dual rotor parameters: the energy-momentum of the particle biases the usual spacetime, effectively ``freezing'' the dual rotor $R_\text{dual}$ to follow $R_\text{usual}$ with a finite stiffness. An inertial frame is modelled as $\Omega=1$, which is not the same as $J=0$.

When the particle is accelerated (non-gravitationally), the usual rotor $R_\text{usual}$ attempts to change instantaneously due to the applied force. However, the dual rotor $R_\text{dual}$, bound by the particle's rest mass, cannot respond immediately — it lags behind. This creates a non-trivial torsional mismatch yielding $\Omega_\text{biv} \neq 0$ and a non-zero contribution to the scalar $J$. The particle experiences this mismatch as inertial resistance — the familiar $F = ma$.

Crucially, the ``cost'' of misalignment is proportional to the rest mass $m$, explaining why heavier particles resist acceleration more strongly. Inertia is therefore not a property of spacetime reacting to mass, but of the particle's own dual spacetime refusing to realign instantaneously.

\subsection{Massless particles}

Massless particles ($m = 0$), such as photons, exhibit no such rigidity. Their usual and dual rotors are perfectly resonant: any change in $R_\text{usual}$ is instantaneously mirrored in $R_\text{dual}$, maintaining $J = 0$ exactly in all frames. Thus, photons experience zero torsional mismatch and zero inertial resistance — they always propagate at $c$ with $\Omega_\text{biv} = 0$. This is why the speed of light is both constant and the universal speed limit: massless particles pay no torsional price for frame changes.

\subsection{The Equivalence Principle as Algebraic Identity}

Consider Einstein's elevator accelerating upward with acceleration $a$. An observer inside feels a fictitious force $ma$ downward. In the dual spacetime picture:

\begin{itemize}
\item The observer (massive) has their usual rotor accelerated by the elevator floor, while their dual rotor lags → torsional mismatch → perceived downward force.
\item A photon entering the elevator bends due to the observer's torsional field, but the photon itself maintains $J = 0$ — it follows a null geodesic in its own dual-flat spacetime.
\end{itemize}

In a gravitational field (elevator at rest on Earth), the floor must accelerate upward relative to free fall to stay stationary, producing the same rotor mismatch as uniform acceleration in flat spacetime.

Thus, the equivalence principle is not an empirical coincidence but an algebraic identity: gravitational and inertial forces are indistinguishable because they are the same dual torsional phenomenon. Free fall is the state where $J = 0$ globally (dual rotors perfectly aligned), which massless particles achieve eternally and massive particles only in the absence of non-gravitational forces.

\subsection{Consequences and Predictions}

\begin{itemize}
\item The arrow of time distinction between usual and dual spacetimes explains why massive particles have a preferred rest frame: the dual time reversal breaks perfect symmetry, generating the mass gap.
\item Inertial mass reduction or elimination becomes possible by externally exciting the dual rotor $\theta$ to synchronize with $\phi$, reducing the effective rigidity — a direct pathway to propellantless propulsion.
\item The theory predicts that sufficiently high-frequency fields coupling to the dual spacetime could make massive objects behave as effectively massless over short timescales, allowing transient superluminal phase velocities (without violating causality, as information remains luminal).
\end{itemize}

Inertia and gravity are unified not by geometry, but by the internal torsional dynamics of particle-intrinsic dual spacetimes. Mass is dual rotor stiffness; motion is the cost of breaking parallelism. With this understanding, both inertia and gravity cease to be fundamental — they become engineerable.

The dream of gravitational and inertial control is no longer science fiction; it is a direct consequence of abandoning the continuum and recognizing spacetime as paired, particle-local, and torsionally active.

\subsection{Upper Limit of Acceleration and the Inertial Reversal Phenomenon}
\label{subsec:acceleration-limit}

The double spacetime theory predicts, as one of its central falsifiable consequences, that acceleration itself possesses an absolute upper limit. The usual-sector boosts are non-compact ($\phi \in \mathbb{R}$) and, taken alone, would permit unbounded rapidity. The dual sector, however, is intrinsically compact: the dual generators $\Gamma_a$ square to $-1$, so the dual rotor $R_\text{dual} = \exp(\theta/2 \, \hat q)$ is governed by the elliptic phase $\theta \bmod 4\pi$. Because inertial response arises from the dual rotor tracking the usual rotor with finite stiffness, the rate at which $\theta$ can be displaced before reaching its periodic boundary sets a strict ceiling on the proper acceleration that a given particle can coherently sustain. We denote this material-intrinsic ceiling by $a_\text{crit}$. Unlike the speed of light, $a_\text{crit}$ is not universal: it depends on the particle species and on the compactification scale of its dual sector, so different atomic constituents possess different critical accelerations.

For proper accelerations $a < a_\text{crit}$, the dual phase $\theta$ tracks $\phi$ smoothly and the familiar Newtonian relation $F = ma$ is recovered. At the instant $a$ approaches $a_\text{crit}$, however, $\theta$ is driven through the compact boundary at $4\pi$ and undergoes a topological wrap-around: the dual rotor re-sutures its phase on the opposite side of the $S^3$-like dual sector. This wrap inverts the sign of the torsional mismatch bivector $\Omega_\text{biv}$, and therefore reverses the direction of the inertial response over a single dual cycle. We refer to this as the inertial reversal (or acceleration reversal) phenomenon: a sudden, discrete flip of the effective $F = ma$ vector triggered when the applied acceleration crosses $a_\text{crit}$. The torsional strain energy that was being stored in the rigid dual sector just below threshold is released abruptly during the wrap, producing an impulsive burst of internal energy in addition to the kinematic reversal itself.

We propose that this mechanism is observationally realised in the anomalous airbursts of large meteoroids. The hypersonic deceleration experienced by bolides entering the Earth's atmosphere can reach values that, for the constituent atoms of the body, are plausibly of order $a_\text{crit}$. Long-standing puzzles in the bolide community---most prominently the Tunguska event (1908) and the Chelyabinsk superbolide (2013), whose disintegration altitudes and total radiated energies are notoriously difficult to reconcile with kinetic-energy dissipation models based on aerodynamic ablation and fragmentation alone---are naturally explained as collective inertial reversals: when the deceleration of the meteoroid crosses the material-specific $a_\text{crit}$, the dual phases of its constituent atoms wrap simultaneously, dumping the accumulated torsional strain energy as an explosive, predominantly radiative pulse. The dependence of $a_\text{crit}$ on composition would also account for the observed variability in airburst altitudes and yields between different meteorite classes (e.g.\ carbonaceous versus ordinary chondrites). A quantitative derivation of $a_\text{crit}$ from the particle-intrinsic compactification scale remains open.

\section{Strong-Field Regime: Layered Gravitational Reversal and the New Black Hole Paradigm}
\label{sec:strongfield}

The double spacetime theory predicts a radically new structure for the endpoints of stellar collapse. The Killing form's sign asymmetry — positive for usual-spacetime boost dominance (attraction) and negative for dual-spacetime rotation dominance (repulsion) — does not simply reverse gravity once and for all. Instead, as density increases inward, the torsional mismatch $\Omega_\text{biv}$ oscillates in dominance between the two sectors, producing alternating layers of attraction and repulsion. This yields a stratified, onion-like interior: attraction → repulsion → attraction → repulsion → ..., with exactly one reversal per period of the interference factor. The amplitude of that alternation is settled in closed form in Section~\ref{subsec:layer-spectrum}: the barriers separating the layers sharpen without bound toward the centre, while the plateaux between them are exponentially screened, so the interior is not a region of ever-stronger force but a stack of ever-sharper walls enclosing ever-quieter interiors.

\subsection{Layered Torsional Structure from the Killing Form}

The scalar $J = \frac12\sum_a(\alpha_a^2-\beta_a^2)$ changes sign each time the dominant contribution flips from $i\Gamma_a$ (boost-like, $+$) to $\Gamma_a$ (rotation-like, $-$). At extreme densities, the dual rotor parameters $\theta$ grow nonlinearly, driving successive resonances that flip the sign of $J$ repeatedly as radius decreases.

The resulting effective potential features:
\begin{itemize}
\item Outer layers: standard attractive gravity ($J > 0$),
\item First reversal: repulsive shell ($J < 0$) that halts infall (core bounce in supernovae),
\item Deeper layers: attractive zones that re-accelerate matter inward, separated from the repulsive shells by sharp barriers,
\item Alternating deeper shells, with barrier strength growing and plateau strength decaying, culminating in chaotic torsional turbulence.
\end{itemize}
On the admissible cone $|J|$ cannot exceed $\Jbound$. Apparent ``unbounded'' growth of the mismatch is therefore a chart artefact: the classical radial rapidity $\varphi\to\infty$ as $r\to r_s^+$ leaves the cone.

\subsection{Quasi-Horizon Cutoff}
\label{subsec:quasi-horizon}

Classical Schwarzschild has $\varphi\to\infty$ and $A\to 0$ as $r\to r_s^+$. Identifying the three-axis admissible ceiling $|J|=\Jbound$ with a radial pure-boost slice $J=\tfrac12\varphi^2$ inverts to the finite cutoff
\[
\varphi_{\max}=\frac{\pi\sqrt{3}}{2},\qquad A_{\min}=\mathrm{e}^{-2\varphi_{\max}}=\mathrm{e}^{-\pi\sqrt{3}}.
\]
The corresponding radius is $r_\star=r_s/(1-A_{\min})$, equivalently $r_\star/r_s=(1-A_{\min})^{-1}$. One has $0<A_{\min}<10^{-2}$ and $1<r_\star/r_s<1.005$. The saturated chart freezes $\varphi=\varphi_{\max}$ and $A=A_{\min}$ for $r\le r_\star$, and uses the classical formulae for $r>r_\star$. At $r=r_\star$ the two descriptions agree. The saturated factor $A$ remains strictly positive: a true one-way horizon $A=0$ does not form. (By contrast the classical factor vanishes at $r=r_s$.)

This cutoff is a working identification of the three-axis ceiling on a single boost axis. A one-axis admissible boost saturates only at $\varphi=\pi/2$, where $J=\pi^2/8=\Jbound/3$. The inverted value $\varphi_{\max}>\pi/2$ therefore lies outside the admissible cone. The construction saturates the three-axis bound, not the single-axis cone.

The dictionary~\eqref{eq:jt-dictionary} converts the rapidity ceiling into a finite envelope for $T$: $\varphi\le\varphi_{\max}$ implies $r^2 T\le 4(\cosh\varphi_{\max}-1)$, and
\[
26<4\bigl(\cosh(\pi\sqrt{3}/2)-1\bigr)<27.
\]
A one-axis admissible boost $\varphi\le\pi/2$ yields the tighter envelope $r^2 T<6.1$. Classical Schwarzschild has $r^2 T\to\infty$ as $r\to r_s^+$; the cutoff therefore bounds $T$ as well as $\varphi$.

The elliptic sector supplies the matching floor. On the admissible cone $\sqrt{2|J|}\le\varphi_{\max}<\pi$, so the cosine branch of $E(J)$ never leaves the range on which $\cos$ is strictly decreasing. Its deepest value is reached at saturation,
\[
-7.7<4\bigl(\cos\varphi_{\max}-1\bigr)<-7.6,
\]
and together with the ceiling this confines the dimensionless teleparallel density to $-8<r^2 T<27$. Repulsive layers have a bounded depth: gravity can reverse, but not with unlimited strength. Both ends of the window are attained, at $J=\pm\Jbound$.

Inside $r\le r_\star$ the redshift floor $A=A_{\min}>0$ is held fixed, so the field seed $J_{\mathrm{field}}$ vanishes. The constant-$A$ coframe still carries angular Weitzenb\"ock torsion, so $T$ itself does not vanish---the dual of a constant boost, which has $J\neq 0$ with $T=0$.

This layered structure prevents both singularities and event horizons. Matter never reaches a central singularity. Instead, the core settles into a quasi-stable, chaotic condensate of extreme-density torsional energy, continuously oscillating between attractive and repulsive phases.

Conventional Kerr-type rotating black holes are ruled out: the theory forbids global frame-dragging of the required magnitude, as angular momentum is redistributed into dual rotor longitudinal modes rather than spacetime rotation.

\subsection{Reinterpretation of Core-Collapse Supernovae and Remnant Formation}

In massive star collapse, the first repulsive layer ($J < 0$) triggers explosive bounce at nuclear densities, powering the supernova without fine-tuned neutrino physics. Deeper attractive layers then recapture portion of the ejecta, forming the stratified remnant. The observed diversity of supernova energies and remnant masses emerges naturally from the number and strength of torsional layers excited.

\subsection{Pulsars as Longitudinal Dual-Torsion Oscillators}

Observed neutron stars are not rapidly rotating magnetized conductors. The double spacetime theory denies both millisecond spin periods and dipolar magnetic fields as the primary energy source.

Instead, pulsar radio pulses arise from ultra-high-frequency longitudinal oscillations of the dual rotors $\theta$ on the star's surface. These torsional waves propagate radially at nearly $c$, modulated by the chaotic interior layering, producing beams via curvature emission in the torsional shock zones.

Key predictions that surpass the rotating magnetar model:
\begin{itemize}
\item The theoretical minimum pulse period is determined by the Planck-scale double rotor frequency, which is orders of magnitude shorter than the observed $\sim 1$ ms limit, allowing us to explain short-period pulsars that could not be explained by the rotating pulsar picture.
\item Pulsar ``glitches'' and period derivatives $\dot{P} < 0$ (period shortening) are the result of the pulsar losing energy causing the vibrational interface to shrink, accelerating the vibrational frequency over time, a natural mechanism not seen in rotation-only models.
\item Magnetar flares and fast radio bursts are violent flips between attractive/repulsive layers, releasing torsional strain without requiring $10^{15}$ G fields.
\item The death line in the $P$-$\dot{P}$ diagram corresponds to the point where interior layering stabilizes and longitudinal modes damp below detectability.
\item The neutron star maximum mass is capped not by quark matter but by the same torsional layering that limits nuclear stability, predicting a strict upper bound $\sim 2.5 M_\odot$ and explaining the absence of pulsars above $\sim 3 M_\odot$.
\end{itemize}

The Crab pulsar's observed spin-down luminosity is reinterpreted as torsional wave leakage rather than magnetic braking, perfectly matching the energetics without invoking unrealistic magnetic field decay.

\subsection{The New Black Hole Paradigm: Torsional Stars}

What observers call ``black holes'' are in fact torsion stars — quasi-stable, layered condensates with no event horizon and no central singularity. The innermost chaotic core is a Planck-density plasma in perpetual torsional turbulence, continuously radiating via rotor tunneling (modified Hawking process). Information is preserved in the dual rotor phases.

Gravitational-wave signals from mergers feature repeated echoes from internal layer reflections, with characteristic frequencies tied to the torsional oscillation spectrum — already hinted in LIGO/Virgo data reanalyses.

The theory thus resolves the information paradox, explains the absence of Kerr spin signatures, and predicts that the most massive ``black holes'' ($\gtrsim 100 M_\odot$) will exhibit periodic modulations in accretion luminosity as inner repulsive layers periodically eject material.

\subsection{Central Floating Core: Zero Time Dilation at the Center}

A common point of confusion in discussions of stellar interiors is the behaviour of gravitational time dilation. In GR, for a spherically symmetric static mass distribution, gravitational time dilation increases monotonically toward the centre: proper time runs most slowly at \(r=0\), where the gravitational potential is deepest.

In the double spacetime theory, the situation is markedly different. Because there is no background continuous spacetime manifold, gravitational effects—including both the effective acceleration and time dilation—arise solely from the torsional mismatch between the particle-intrinsic usual and dual rotors.

At the exact geometric centre of a spherically symmetric torsion star, perfect symmetry forces the usual-spacetime boost parameters \(\phi_a\) induced by mass elements in opposite directions to cancel completely. The residual mismatch angle \(\delta\phi_a = \phi_a - \theta_a\) therefore vanishes at \(r=0\), yielding \(J=0\) and a complete restoration of torsional balance.

Consequently:

\begin{itemize}
  \item The effective gravitational acceleration is zero at the center (in Newtonian mechanics, GR, and DST).
  \item Gravitational time dilation is also zero at the centre: proper time at \(r=0\) runs at the same rate as in the vacuum far from the star.
\end{itemize}

This vanishing of time dilation at the core is a distinctive, observable-in-principle prediction of the double spacetime theory. It contrasts sharply with GR, where the central region experiences the strongest time dilation.

The zero-dilation central core also contributes to the stability of torsion stars by eliminating the infinite central compression and extreme clock-rate gradients that lead to singularities in the classical GR description.

\subsection{Implications for Planetary and Stellar Cores: Instability and Fluctuations}

The floating core paradigm, where $J(r=0) \approx 0$ and effective gravity vanishes at the center, implies that the cores of planets and stars are not inert, stable regions but dynamically active zones prone to torsional instability. Unlike the monotonic increase in gravitational potential assumed in continuum models, the double spacetime theory predicts that the central floating core is a metastable state, susceptible to perturbations that trigger oscillatory transitions between attractive and repulsive torsional layers. These fluctuations arise from the delicate balance of dual rotor phases in the isotropic core, where even minor density variations or external influences can shift the resonance condition $\beta(0) \to 1$, leading to transient excursions of $J$ away from zero. Such excursions manifest as localized bursts of torsional energy, effectively generating mini torsion stars within the core. These mini torsion stars can induce seismic activity, magnetic field fluctuations, and energy transport anomalies observable at the surface. In stellar interiors, these core fluctuations may contribute to phenomena such as solar flares, sunspots, and variability in luminosity, as the torsional energy propagates outward through the layered structure. In planetary contexts, core torsional instability could explain geomagnetic reversals, mantle plumes, and episodic volcanic activity, as the torsional dynamics couple to convective processes in the overlying layers. Thus, the floating core is not a quiescent region but a crucible of dynamic torsional phenomena with far-reaching implications for planetary and stellar behavior.

\section{Time Dilation in Double Spacetime Theory: The Dynamical Role of Dual-Spacetime Proper Time}
\label{sec:timedilation}

Time dilation in the double spacetime theory emerges from the interplay between the usual spacetime and the dual spacetime, each endowed with its own proper time. The usual spacetime proper time $\tau$ governs the familiar kinematics observed in macroscopic Minkowski space, while the dual spacetime proper time $\tilde{\tau}$, encoded in the compact rotor phases $\theta$, drives the internal torsional dynamics that manifest as gravitational and inertial effects.

\subsection{Derivation of the Effective Lorentz Factor}

The effective Lorentz factor $\gamma_\text{eff} = dt/d\tau$ is derived directly from the algebraic action of the total rotor
\[
R_\text{total} = R_\text{usual} R_\text{dual},
\]
where
\[
R_\text{usual} = \exp\!\left(\frac{\phi}{2} \hat{p}\right) = \cosh\frac{\phi}{2} + \hat{p}\: \sinh\frac{\phi}{2}, \quad \hat{p}^2 = +1,
\]
generates boosts in the usual spacetime, and
\[
R_\text{dual} = \exp\!\left(\frac{\theta}{2} \hat{q}\right) = \cos\frac{\theta}{2} + \hat{q} \sin\frac{\theta}{2}, \quad \hat{q}^2 = -1,
\]
generates rotations in the dual spacetime.

To reveal the origin of the interference term explicitly, consider the pure time vector $X = ct \, j$ ($j^2 = -1$) in the usual spacetime basis. The transformed vector is $\tilde{X} = R_\text{total} \, X \, R_\text{total}^\dagger$.

Due to the time-reversal duality, the inverse dual map from dual to usual spacetime is $X \mapsto X(-i)$. This map naturally yields the alignment $\hat{q} (-i) = -\hat{p}$ in the vacuum limit, transforming the dual rotor into an effective opposite-oriented rotor
\[
R_\text{dual}^\text{eff} = \cos\frac{\theta}{2} - \hat{p} \sin\frac{\theta}{2}.
\]
The negative sign in the sin term arises directly from this inverse mapping, reflecting the anti-synchronization required for torsion-free flat spacetime.

The total effective rotor is then
\[
R_\text{eff} = R_\text{usual} R_\text{dual}^\text{eff} = \left( \cosh\frac{\phi}{2} + \hat{p} \sinh\frac{\phi}{2} \right) \left( \cos\frac{\theta}{2} - \hat{p} \sin\frac{\theta}{2} \right).
\]

Expanding gives the scalar part
\[
s = \cosh\frac{\phi}{2} \cos\frac{\theta}{2} - \sinh\frac{\phi}{2} \sin\frac{\theta}{2}
\]
and the bivector coefficient
\[
b = -\cosh\frac{\phi}{2} \sin\frac{\theta}{2} + \sinh\frac{\phi}{2} \cos\frac{\theta}{2}.
\]

The net scaling of the time component arises from the squared magnitude of $R_\text{eff}$ acting on the time vector via the sandwich product. For a rotor $R_\text{eff} = s + b\hat{p}$ with $\hat{p}^2 = +1$ and $\hat{p}$ anticommuting with the time generator $j$, the sandwich action on $X = ct\, j$ yields a time component scaled by $s^2 + b^2$, defining the effective Lorentz factor unambiguously as
\[
\gamma_\text{eff} \;\equiv\; s^2 + b^2.
\]
It is convenient also to introduce the scalar interference factor
\[
\gamma_s \;\equiv\; s \;=\; \cosh\frac{\phi}{2} \cos\frac{\theta}{2} - \sinh\frac{\phi}{2} \sin\frac{\theta}{2},
\]
which, unlike $\gamma_\text{eff}$, may change sign and will serve as the order parameter distinguishing attractive ($\gamma_s > 0$) from repulsive ($\gamma_s < 0$) torsional regimes below. The negative sign in the interference term is fixed by the time-reversed duality of the dual spacetime.

A short calculation collapses the squares:
\begin{equation}
\label{eq:gamma-eff}
\gamma_{\mathrm{eff}}(\alpha,\beta)=\cosh\alpha-\sinh\alpha\sin\beta.
\end{equation}
This factor is always at least $e^{-|\alpha|}$, hence never zero: clocks slow, they do not stop. If the dual sector is unexcited, $\beta=0$, one recovers the special-relativistic factor $\cosh\alpha$. If there is no usual boost, $\alpha=0$, then $\gamma_{\mathrm{eff}}=1$. The balanced locus $\alpha=\beta$, where $J=0$ on one axis, is a different kinematics: $\gamma_{\mathrm{eff}}(\alpha,\alpha)\neq\cosh\alpha$ as soon as $\alpha\neq 0$. Special relativity is the unexcited dual sector $\theta=0$, not the vanishing of $J$.

\subsection{Vacuum Limit and Massless Particles}

Torsion-free kinematics is $\theta=0$ (equivalently $R_{\mathrm{dual}}=1$ on a pure usual boost), which yields $\gamma_{\mathrm{eff}}=\cosh\phi$. Balanced massive states have $J=0$ with $\theta=\phi\neq 0$ and a different factor $\gamma_{\mathrm{eff}}=\cosh\phi-\sinh\phi\sin\phi$. Photons, in this reading, maintain $\theta=0$ at all velocities: they pay no dual-rotor price for a change of frame.

\subsection{Massive Particles and Gravitational Effects}

For massive particles, rest mass is the rigidity that causes the dual rotor to lag behind the usual rotor under acceleration. The resulting deviation $\delta\theta = \theta - \phi$ modulates $\gamma_{\mathrm{eff}}$. We read that modulation as gravitational redshift.

\subsection{Sign Inversion of the Scalar Interference Factor and Repulsive Layers}

When dual-spacetime rotations dominate ($\theta/2$ sufficiently large), the sinusoidal term can overpower the hyperbolic growth:
\[
\tan\left(\frac{\theta}{2}\right) > \coth\left(\frac{\phi}{2}\right).
\]
Then $\gamma_s$ becomes negative, while $\gamma_\text{eff}=s^2+b^2$ remains strictly positive. The sign of $\gamma_s$ is the order parameter of repulsive layers. It is not a reversal of causality: proper time never runs backward.

This $\theta$-centric formulation makes manifest that time dilation is fundamentally a torsional phenomenon originating in the particle-intrinsic dual spacetime. The usual rapidity $\phi$ provides the kinematic baseline, while deviations in the compact dual angle $\theta$ encode the gravitational and inertial corrections that distinguish GR from flat-spacetime kinematics. External coherent excitation of $\theta$ thus offers a pathway to manipulate time dilation and effective gravitational fields.

\section{Scale-Invariant Torsional Layers: The Strong Force as Gravitational Repulsion and Nuclei as Primordial Torsion Stars}
\label{sec:nuclear}

The layered reversal is scale-invariant in the dimensionless variable \(x=\ell/(2r)\). Under the equal-scale ansatz \(\alpha(r)=\beta(r)=\ell/r\) the entire sequence depends on \(x\) alone; femtometer and MeV scales enter only through a characteristic length \(\ell\) and the particle masses. The same mechanism that stratifies a collapsing star is therefore available at nuclear densities. We propose that this is the strong force: gravity in its layered phase, not a second coupling. The Newtonian gravity fine-structure constant \(\alpha_G=G m_p^2/\hbar c\sim 10^{-38}\) must not be conflated with an \(O(1)\) nuclear strength. Near a torsional node, \(1/\gamma_s\) diverges; that is a kinematic amplification of a channel already present, not a derivation of the strong coupling from \(G\).

One structural point must be settled at the outset. The equal-scale ansatz sets \(\alpha_a=\beta_a\) on every axis, so \(J\) vanishes identically while \(M\) stays positive whenever some rapidity is nonzero. Nuclear matter, in this reading, sits on the balanced locus of Section~\ref{sec:gravity}: massive, with \(\Omega\neq 1\), and with \(J=0\). What alternates across the layers is not \(J\) but the interference factor \(\gamma_s\) of Section~\ref{sec:timedilation}. Attraction and repulsion below are the sign of \(\gamma_s\). \(J\) and \(\gamma_s\) are independent order parameters; nuclear layering belongs to the second.

\subsection{The Closed-Form Layer Spectrum}
\label{subsec:layer-spectrum}

Written in the dimensionless variable \(x=\ell/(2r)\), the equal-scale interference factor is
\[
\gamma_s(x)=\cosh x\cos x-\sinh x\sin x .
\]
Its derivative collapses to a single product,
\begin{equation}
\label{eq:gammaS-derivative}
\gamma_s'(x)=-2\cosh x\,\sin x ,
\end{equation}
so the entire layer structure is governed by the zeros of \(\sin\). The critical points of \(\gamma_s\) are exactly the integer multiples of \(\pi\), and the extremal values are
\[
\gamma_s(n\pi)=(-1)^n\cosh(n\pi),\qquad n=0,1,2,\dots
\]
On each interval \([n\pi,(n+1)\pi]\) the factor is strictly monotone---decreasing for even \(n\), increasing for odd \(n\)---so each interval carries exactly one torsional node, and the sign of \(\gamma_s\) strictly alternates from one plateau to the next. Plateaux never vanish; only finitely many nodes lie in any bounded range of \(x\).

Ordered outward in radius the pattern is: attraction for \(r>r_1\); a barrier at \(r_1\); repulsion on \((r_2,r_1)\); a barrier at \(r_2\); attraction on \((r_3,r_2)\); and so on. The onion of Section~\ref{sec:strongfield} is not an extra assumption. It is this derivative, one barrier per \(\pi\)-interval of \(x\), with no two consecutive layers of the same sign.

A node forces \(\tan x=1/\tanh x>1\), so the \(n\)-th node sits in
\[
n\pi+\frac{\pi}{4}<x_n<n\pi+\frac{\pi}{2},
\]
or in radius
\[
\frac{\ell}{(2n+1)\pi}<r_n<\frac{2\ell}{(4n+1)\pi}.
\]
Thus \(r_n=\Theta(1/n)\), every node obeys \(r_n<2\ell/\pi\), and the outermost node is the largest layer radius. The first two nodes satisfy
\[
\frac16<\frac{r_2}{r_1}<\frac{4}{5\pi}\approx 0.255.
\]

The plateau amplitudes grow at least geometrically, \((\cosh\pi)^n\le\cosh(n\pi)\) with \(\cosh\pi>11\), while never exceeding \(\mathrm{e}^{n\pi}\). Successive layers amplify by an ever larger factor that stays strictly below \(\mathrm{e}^\pi\approx 23.14\). Because \(\gamma_s\) sits in the denominator of the effective force \(F=k/(\gamma_s r^2)\), the mid-layer amplitude at the \(n\)-th plateau \(r=\ell/(2n\pi)\) is
\[
|F|=\frac{k}{\ell^2}\cdot\frac{4\pi^2n^2}{\cosh(n\pi)},
\]
which decreases in \(n\) for \(n\ge 1\). What grows inward without bound is the barrier structure---the nodes at which \(\gamma_s\) vanishes and \(1/\gamma_s\) diverges. The plateaux between barriers become exponentially quiet, successive factors falling like \(\mathrm{e}^{-n\pi}\). In radius the screening is \(1/\cosh\bigl(\ell/(2r)\bigr)\): the deep interior of a torsional condensate is quasi-free between barriers and confined at them. That is the qualitative pattern of nuclear confinement, read off the algebra rather than postulated.

A barrier is a force divergence, not a clock freeze: at a node \(\gamma_{\mathrm{eff}}=\gamma_b^2>0\) by~\eqref{eq:gamma-eff}. The ratios \(r_{n+1}/r_n\) depend only on the dimensionless roots \(x_n\), so they are identical at nuclear and at stellar scales; \(\ell\) fixes the absolute radii and cancels from every ratio. That is the content of the scale invariance in the title of this section, and the only sense in which nuclear and stellar layering are the same structure.

\subsection{Torsional Layers at Nuclear Densities}

As baryon density approaches nuclear saturation
\(\rho_0\sim (2\text{--}3)\times 10^{14}\,\mathrm{g/cm^3}\)
(equivalently \(n_0\sim 0.16\,\mathrm{fm}^{-3}\)),
the dual rotor parameters \(\theta\) dominate the mismatch. Densities of order \(10^{18}\,\mathrm{g/cm^3}\) are not ordinary nuclear interiors. The sign flips of \(\gamma_s\) are organised on femtometer scales set by \(\ell\). Taking \(\ell\sim\lambda_\pi\approx 1.41\,\mathrm{fm}\), the charged-pion Compton wavelength, is a physical identification, not a derivation. Because \(\lambda_\pi<\pi/2\,\mathrm{fm}\), every equal-scale node then lies below \(1\,\mathrm{fm}\), and the outermost node is confined to
\[
0.44\,\mathrm{fm}<r_1<0.91\,\mathrm{fm}.
\]
That is the empirical hard-core region near \(0.5\,\mathrm{fm}\). A repulsive shell at \(1\)--\(2\,\mathrm{fm}\) would require a longer \(\ell\).

The layer sequence at nuclear densities then reads:

\begin{itemize}
\item Outermost barrier: the first node \(\gamma_s=0\), identified with the short-range repulsion of nucleon--nucleon scattering, at \(0.44\)--\(0.91\,\mathrm{fm}\) if \(\ell=\lambda_\pi\).
\item Repulsive shell inside it (\(\gamma_s<0\)), extending inward to the second node, whose radius is between one sixth and one quarter of the first.
\item Intermediate attractive plateau (\(\gamma_s>0\)). The observed binding energy per nucleon is \(\mathrm{BE}/A\sim 8\,\mathrm{MeV}\) near iron; an optical-model well depth of order \(40\)--\(50\,\mathrm{MeV}\) is a distinct input.
\item Deeper layers: strictly alternating attraction and repulsion, one barrier per \(\pi\)-interval of \(x\), with exponentially quieter plateaux. The barriers halt further collapse and are the geometric counterpart of quark confinement; the quiet plateaux between them are a nearly free interior.
\end{itemize}

Pauli's exclusion principle is no longer fundamental: it is dynamically generated by the first repulsive torsional layer. Fermions cannot occupy the same state because the dual rotor $\theta$ phases enter a repulsive regime when two identical particles attempt spatial overlap, producing an effective fermi pressure without invoking quantum statistics a priori. Bosons, lacking this phase rigidity, experience only the attractive layers until much higher densities.

In short: we propose that the strong force is gravity in its strong-field, layered phase. QCD color charge is then a misreading of torsional charge carried by dual-rotor phases; we do not identify \(\alpha_s\) with the discrete lattice floor \(\varepsilon_N\).

\subsection{Atomic Nuclei as Primordial Torsion Stars}

Every atomic nucleus is a microscopic analog of the torsion stars described in Section~\ref{sec:strongfield}, and by the previous subsection the two share not merely a qualitative resemblance but the same dimensionless node sequence. Reading inward from the nuclear surface, the derived structure is:

\begin{itemize}
\item Outermost repulsive barrier at \(r_1\): generates the hard-core repulsion in NN scattering and the nuclear surface tension,
\item Repulsive shell just inside it, terminating at \(r_2<r_1/4\),
\item Attractive plateau between \(r_3\) and \(r_2\): mediates binding between nucleons,
\item Alternating deeper barriers and plateaux, the barriers ever sharper and the plateaux exponentially quieter, so that the innermost region is nearly interaction-free between confining walls.
\end{itemize}

The amplification is entirely in the barriers; the intervening interior is screened by \(1/\cosh\bigl(\ell/(2r)\bigr)\).

Saturation of nuclear density is the one nuclear regularity the layer picture claims to explain, and it is also where the standard empirical constants disagree. Constant nucleon number density is exactly \(R=r_0A^{1/3}\): a uniform sphere of that radius has density \(3/(4\pi r_0^3)\) for every \(A\). Numerically \(r_0=1.2\,\mathrm{fm}\) gives
\[
0.13\,\mathrm{fm}^{-3}<\frac{3}{4\pi r_0^3}<0.14\,\mathrm{fm}^{-3},
\]
which is not the empirical saturation value \(n_0=0.16\,\mathrm{fm}^{-3}\). The coefficient consistent with \(n_0\) is \(r_0^\star\approx 1.14\,\mathrm{fm}\). The two commonly quoted nuclear constants disagree at the ten-per-cent level; a layer identification must choose between them.

Charge independence, spin--orbit coupling, and magic numbers are read as resonance conditions in the dual-rotor spectrum. Discrete finiteness of the rotor image at fixed compactification \(N\) yields only \(A_{\max}(N)<\infty\). The layer geometry itself yields no such cutoff: if \(\ell\propto A^{1/3}\), the ratio of the nuclear surface to any layer radius is independent of \(A\); if \(\ell=\lambda_\pi\) is \(A\)-independent, already at \(A=1\) the surface \(r_0=1.2\,\mathrm{fm}\) lies outside every node. The estimate \(A\sim 300\) remains an external heuristic.

Nuclear fission and fusion are torsional layer transitions: fission occurs when an excited dual mode drives \(\gamma_s\) through a node from positive to negative, explosively repelling fragments; fusion succeeds when two nuclei tunnel through the outermost barrier into the attractive plateau behind it.

\subsection{Experimental Signatures and Immediate Predictions}

\begin{itemize}
\item Ultra-high-energy nucleus-nucleus collisions (LHC, RHIC) should produce transient \(\gamma_s<0\) fireballs that decay via explosive hadron jets rather than hydrodynamic flow — already hinted in anomalous flow data.
\item Precision measurements of the nuclear equation of state above saturation density will reveal oscillatory behavior in pressure vs. density, with characteristic repulsive peaks at $\sim$ 2–4 $\rho_0$.
\item The ratio of successive barrier radii is a pure number, confined to \((1/6,\,4/(5\pi))\) for the first pair and independent of \(\ell\). The same ratio must therefore appear in the NN repulsive core and in the layer spacing of compact remnants; a measured mismatch between the two would falsify the scale-invariance claim without any appeal to \(\ell\).
\item Because the plateaux are screened by \(1/\cosh\bigl(\ell/(2r)\bigr)\), probes that resolve well inside the outermost barrier should see an interaction that weakens exponentially with decreasing distance, interrupted only at the discrete barrier radii.
\end{itemize}

The four interactions are not collapsed into a single force. They are two channels of one geometry. Gravity and electromagnetism share the layered reversal of \(\gamma_s\) and differ only in which phase of the dual rotor sources the mismatch: the kinematic pair $(\phi, \theta)$ tied to inertial mass, or an independent pair $(\alpha, \beta)$ tied to charge. The strong force is the layered phase of the gravitational channel at nuclear densities; electron shells, in the next section, are the analogous layered phase of the electromagnetic channel. Faraday coefficients occupy the same six-generator space as the dual rotor (Section~\ref{sec:unification}). Identifying commuting and anticommuting triples with charged and neutral currents, or recovering a Weinberg angle, remains open.

There is only one geometry — the layered dance of attraction and repulsion in the dual rotor, played simultaneously in two phase channels. Atomic nuclei are the primordial torsion stars of the gravitational channel; atomic electron shells, as we shall see next, are the primordial torsion stars of the electromagnetic channel — and from this single dual-rotor geometry, expressed in two angular channels, all structure emerges.

\section{Time Reversal and Electron Shells via Coulombic Torsional Layers}
\label{sec:electronshells}

In the double spacetime framework, the Coulomb interaction between charged particles can be reinterpreted as arising from torsional mismatch in the dual spacetime sector, analogous to the gravitational case but with electromagnetic origin. Replacing the gravitational constant \(G\) and masses \(M,m\) with the Coulomb constant \(k = 1/(4\pi\epsilon_0)\) and charges (e.g., \(+e\) for the proton and \(-e\) for the electron), the effective force takes the form
\[
F(r) = \frac{k e^2}{\gamma_s(r)\, r^2},
\]
where \(\gamma_s\) is the scalar interference factor introduced in Sec.~\ref{sec:timedilation}, which retains the same algebraic structure from the dual rotor interplay (and which can change sign, unlike the proper-time Lorentz factor \(\gamma_\text{eff} \geq 0\)):
\[
\gamma_s(\alpha(r),\beta(r)) = \cosh\!\left(\frac{\alpha(r)}{2}\right) \cos\!\left(\frac{\beta(r)}{2}\right) - \sinh\!\left(\frac{\alpha(r)}{2}\right) \sin\!\left(\frac{\beta(r)}{2}\right).
\]

Here, \(\alpha(r)\) and \(\beta(r)\) are interpreted as electromagnetic rapidity and dual rotation angle parameters that scale inversely with distance: \(\alpha(r) \propto 1/r\) and \(\beta(r) = \lambda / r\), with \(\lambda\) a characteristic length scale set by the compact dual spacetime.

For the hydrogen atom (\(Z=1\)), this yields
\[
F(r) = \frac{k e^2}{r^2 \left[ \cosh\!\left(\frac{\alpha(r)}{2}\right) \cos\!\left(\frac{\beta(r)}{2}\right) - \sinh\!\left(\frac{\alpha(r)}{2}\right) \sin\!\left(\frac{\beta(r)}{2}\right) \right]}.
\]

The hyperbolic-trigonometric form of \(\gamma_s\) introduces extreme oscillatory behavior. As \(r\) decreases:

\begin{itemize}
  \item At large \(r\), \(\alpha(r),\beta(r) \to 0\), so \(\gamma_s \to 1\), recovering the standard Coulomb attraction.
  \item At smaller radii, the interference term generates points where \(\gamma_s = 0\) (infinite repulsion) and regions where \(\gamma_s < 0\) (force reversal to repulsion).
  \item These discrete radii \(r_n\) satisfy the resonance condition
  \[
  \tanh\!\left(\frac{\alpha(r_n)}{2}\right) \tan\!\left(\frac{\beta(r_n)}{2}\right) = 1.
  \]
\end{itemize}

The force is undefined at a node. A node is a force divergence, not a clock freeze: \(\gamma_{\mathrm{eff}}=\gamma_b^2>0\) by~\eqref{eq:gamma-eff}. Infinite barriers separate neighbouring wells, but they do not by themselves supply a circular orbit.

\subsection{First Coulombic node and characteristic length}

Under \(\alpha(r)=\beta(r)=\lambda/r\), the first positive node of \(\gamma_s\) lies in \(\pi/4<x_1<1\) with \(x=\lambda/(2r)\). If that outermost node is identified with the Bohr radius \(a_0\), the characteristic length is confined to \(\frac\pi2 a_0<\lambda<2a_0\). The identification \(r_1=a_0\) is a proposal; \(a_0\) itself is an external input.

\subsection{Coulombic circular orbits}

The centripetal balance \(mv^2/r=ke^2/(\gamma_s r^2)\) rearranges to
\[
v^2=\frac{ke^2}{m\,\gamma_s\,r}.
\]
Circular speed is real only where \(\gamma_s>0\). On a repulsive layer the circular \(v^2\) is negative. In the Coulomb limit \(\gamma_s=1\) one recovers \(v^2=ke^2/(mr)\). Bound circular orbits therefore exist only in attractive layers---the exterior Coulomb well \(r>r_1\) and the inward attractive strata between later barriers.

Under the equal-scale ansatz the resonance condition is \(\tanh x\cdot\tan x=1\), the same as \(\gamma_s=0\), so the layer spectrum of Section~\ref{subsec:layer-spectrum} applies verbatim: radii scale as \(r_n=\Theta(1/n)\), not as \(\Theta(n^2)\). The first two nodes satisfy \(r_2/r_1<1\), while Bohr bookkeeping \(r_n=n^2 a_0\) gives \(r_2/r_1=4\). The equal-scale tower runs inward; Bohr shells run outward. Matching the Rydberg series therefore requires a separate radius law such as \(r_n=n^2 a_0\). The barriers are the geometry; the hydrogen spectrum is additional bookkeeping.

\subsection{\texorpdfstring{$Z$-contraction}{Z-contraction}}

On the equal-scale locus, \(\ell\mapsto\ell/Z\) contracts every layer radius by \(Z\), while \(\ell\mapsto Z\ell\) expands them. The substitution \(\alpha,\beta\propto Z/r\) at fixed \(\ell\) is the expanding map: it evaluates the \(Z=1\) profile at \(r/Z\) and moves nodes outward. Observed shrinkage of electron shells with nuclear charge is the identification \(\ell(Z)=\ell/Z\).

The shared discrete floor \(\varepsilon_N=16/(3N^2)\) may be paired with the electromagnetic coupling as \(\alpha_{\mathrm{hyp}}(N)=\varepsilon_N\). Matching CODATA \(\alpha\) then forces \(729<N_*^2<732\), so \(N_*\) lies between \(27\) and \(28\), and \(\alpha_{\mathrm{hyp}}(27)=16/2187\) overshoots \(\alpha\) slightly. This is a diagnostic, not a derivation of \(\alpha\). It is a different channel from the gravitational height of the electron, which yields \(N_*^2\sim 10^{45}\).

\subsection{Energy Conservation in the Dual Structure}

The apparent non-conservative behavior of the manifest potential \(V_\text{eff}(r) = -\int F(r)\, dr\)—with its violent oscillations and sign changes—is resolved by distinguishing between the manifest potential observed in the usual spacetime and the true potential stored in the compact dual-spacetime rotor phases \(\theta_a\).

Total energy is conserved across both sectors:
\[
E_\text{total} = E_\text{kinetic (usual)} + V_\text{eff}(r) + V_\text{true}(\theta) = \text{constant}.
\]
The true potential \(V_\text{true}\) is proportional to the torsional mismatch scalar \(J \propto (\phi - \theta)^2\), acting as a hidden reservoir. During layer transitions, a decrease in \(V_\text{eff}\) (attractive to repulsive) is exactly compensated by an increase in \(V_\text{true}\) as the dual rotor winds, and vice versa.

Photon emission corresponds to the release of stored dual-phase energy when the rotor unwinds into a lower torsional state. This dual bookkeeping ensures strict energy conservation without violating the discrete, particle-local nature of spacetime. Observed spectral differences require the separate Bohr bookkeeping of shell radii discussed above; they are not implied by equal-scale \(\gamma_s=0\) alone.

The Coulombic force therefore accounts for discrete torsional barriers and for circular orbits in attractive layers, without a wave-function postulate. Identifying those barriers with atomic electron shells is a proposal: the layer spectrum is the geometry, its absolute scale and the Bohr bookkeeping are inputs.

\section{Electromagnetism and the Weak Channel in Dual Spacetime}
\label{sec:unification}

Electromagnetism and the weak channel occupy the same six-generator space as the dual rotor. They are not the same map. The dual map $X\mapsto Xi$ is Hodge duality on that space: it exchanges the usual (electric, hyperbolic) and dual (magnetic, cyclic) summands, has period 4, and sends $J\mapsto -J$. Chirality, by contrast, requires a generator that squares to $+1$. The $\mathrm{Cl}(3,1)$ pseudoscalar satisfies $i^2=-1$, so $(1\pm i)/2$ are not idempotent and cannot serve as left/right projectors. The complementary pair built from the spatial axis $e_1$ with $e_1^2=+1$,
\[
P_L = \frac{1 - e_1}{2},\qquad P_R = \frac{1 + e_1}{2},
\]
is idempotent, sums to the unit, and multiplies to zero. Duality sends $e_1$ to the trivector $-e_0 e_2 e_3$, which is not a real multiple of $e_1$: the chirality axis is not a dual-invariant line. The product of the commuting square-$+1$ pair $e_1$ and $e_0e_2$ generates a nonzero left ideal; we do not claim that ideal as an irreducible spinor module.

The Faraday coefficients $(E_a,B_a)$ occupy the same six-generator space as the dual-rotor bivector $\Omega_{\mathrm{biv}}$:
\begin{equation}
\label{eq:faraday-split}
F = F_{\mathrm{usual}}+F_{\mathrm{dual}},\qquad
F_{\mathrm{usual}}=\sum_a E_a\,i\Gamma_a,\qquad
F_{\mathrm{dual}}=\sum_a B_a\,\Gamma_a.
\end{equation}
Electric field is usual-sector, boost-like, noncompact. Magnetic field is dual-sector, rotation-like, compact. Duality $X\mapsto Xi$ acts on coefficients by $(E,B)\mapsto(B,-E)$. On this identification the mismatch scalar is the Faraday quadratic
\[
J(F)=\tfrac12(E^2-B^2),\qquad M(F)=\tfrac12(E^2+B^2).
\]
A pure electric configuration has $J\ge 0$; a pure magnetic configuration has $J\le 0$. That is the Killing-form signature of the two sectors, written in laboratory fields. Fields with $E^2=B^2$ sit on the balanced massive locus $J=0$, $M>0$, not on the vacuum $M=0$.

A circularly polarised wave of helicity $\sigma=\pm 1$ (units $c=1$),
\[
\mathbf{E}=E_0(\hat{x}\cos\psi+\sigma\hat{y}\sin\psi),\qquad
\mathbf{B}=E_0(-\hat{x}\sin\psi+\sigma\hat{y}\cos\psi),
\]
is null at every phase $\psi=kz-\omega t$: $J=0$, $E\cdot B=0$, and $M=E_0^2$. Its signed Poynting component along the beam is $\sigma E_0^2$ at every phase. Laboratory time reversal $(E,B)\mapsto(E,-B)$ reverses that sign; duality preserves it. The two local Faraday invariants $J$ and $E\cdot B$ vanish identically on the wave, so a helicity-odd scalar cannot be a polynomial in those invariants. Superposing the wave onto a background torsion shifts $J$ by a first harmonic in the phase; the four-phase mean of $J$ is the background value. Linear superposition does not displace mean mismatch.

Three maps act on this six-space and must not be conflated. Duality has period 4 and sends $J\mapsto -J$. The usual--dual parameter swap has period 2, also sends $J\mapsto -J$, and coincides with duality only on a pure magnetic field. Laboratory time reversal preserves $J$ and flips $E\cdot B$. Same-axis electric and magnetic generators commute, so a simultaneous rotation of the pair $(E,B)$ is well-defined:
\[
E_a(\omega)=E_a\cos\omega+B_a\sin\omega,\qquad
B_a(\omega)=-E_a\sin\omega+B_a\cos\omega,
\]
with $J(\omega)=J\cos 2\omega+(E\cdot B)\sin 2\omega$. Duality is the special angle $\omega=\pi/2$.

Relative to $e_1$ the six generators split into a commuting triple $(B_1,E_2,E_3)$ and an anticommuting triple $(E_1,B_2,B_3)$. The sandwich $P_R X P_R$ keeps the former and kills the latter; the selection is axis-dependent, not a Lorentz-scalar Weyl projection. A circular wave is the phase quadrature of those two triples, so neither projector selects the helicity sign $\sigma=\pm 1$. Axis chirality is not photon helicity.

The first-order sandwich increment of a rotor $R=1+\varepsilon\Omega$ is the commutator $\Omega X-X\Omega$, not the outer product $F\wedge X$. On a rest particle a pure electric generator produces a spatial kick and every magnetic generator produces none. On $X=e_0+\mathbf{v}$ the commutator is the 4-vector
\begin{equation}
\label{eq:faraday-lorentz}
FX-XF=(E\cdot v)\,e_0+(E-v\times B).
\end{equation}
That is the Lorentz 4-force of the laboratory time-reversed field $(E,-B)$. The familiar 3-force $q(E+v\times B)$ differs from the sandwich by the sign of $B$; laboratory time reversal converts one into the other. A rest frame receives the electric kick $E$ with vanishing power. A velocity along a circular beam yields a transverse kick with no longitudinal component. Maxwell's equations are not derived: the identity is kinematic on the sandwich.

The electromagnetic and weak channels therefore live in one six-space. The dual map exchanges the two sectors and is not a chirality projector. The sign of $J$ distinguishes the noncompact electric sector from the compact magnetic sector. The sandwich on a particle's 4-velocity is the Lorentz force of $(E,-B)$. Minimal coupling $\theta_a\mapsto\theta_a+q A_a$, an extended rotor built from $\int F\,ds$, a Weinberg angle from the mix $\omega$, and a Higgs reading of rotor rigidity remain open. The dual map does not imply a V--A coupling.

\section{Dark Matter-Free Galactic Rotation Curves from Dual-Spacetime Compactness}
\label{sec:darkmatter}

The same dual-rotor geometry that layers nuclei and electron shells unfolds, at cosmological scales, as compactness. We read the gravitational channel, driven by $(\phi,\theta)$ tied to inertial mass, as compactifying the dual sector into galaxy-scale $S^3$ patches; the electromagnetic channel, driven by $(\alpha,\beta)$ tied to charge, unfolds over substantially larger, near-cosmic scales. On the three-sphere chart itself the geometry is definite: the cotangent potential obeys Gauss's law; circular speed is $v^2=(GM/R)\,f(r/R)$; the shape $f$ has a unique minimum on the attractive interval; and attraction ends at $r=\pi R/2$, strictly before $r=2R$. If galaxies are such patches, the systematic distortion of the compact-to-flat map is what observers interpret as missing gravitational mass, and flat rotation curves emerge from baryons alone. That last step is the proposed reading.

\subsection{The Compact Dual Spacetime as 3-sphere}

The dual rotor $R_{\text{dual}} \in \mathrm{SU}(2) \cong S^3$ parameterizes the dual spacetime degrees of freedom intrinsic to every massive particle. The three independent rotation angles $(\theta_1, \theta_2, \theta_3)$ define coordinates on a 3-sphere of compactification radius $R$, while the usual spacetime remains flat and non-compact.

The logarithmic (inverse exponential) map
\[
\log : S^3 \to \mathfrak{su}(2) \cong \mathbb{R}^3
\]
identifies each dual rotor with a point in flat $\mathbb{R}^3$, mapping the 3-sphere onto a ball of radius $\pi R$. This map is the three-dimensional generalization of the azimuthal equidistant projection familiar from cartography.

Just as the azimuthal equidistant projection of $S^2$ centered on the North Pole faithfully represents distances near the pole but distorts areas catastrophically near the South Pole --- mapping a single antipodal point to the entire circumference of the projection boundary --- the logarithmic map of $S^3$ introduces a systematic volume distortion quantified by the Jacobian
\[
\mathcal{J}(\chi) = \frac{\sin^2 \chi}{\chi^2},
\]
where $\chi = r/R$ is the geodesic angle corresponding to the physical distance $r$ in the usual spacetime. This Jacobian equals unity at the origin ($\chi = 0$, faithful representation) and vanishes at the antipode ($\chi = \pi$, maximal distortion). All standard gravitational calculations implicitly assume $\mathcal{J} = 1$; the double spacetime theory reveals that this assumption fails at galactic scales.

\subsection{Gravitational Potential on 3-sphere: The Cotangent Potential}

The torsional mismatch scalar $J$, from which gravity arises, is computed on the compact dual spacetime. The effective gravitational potential for a point mass $M$ is governed by the Green's function of the Laplacian on $S^3$. For the spherically symmetric case, the unique radially harmonic function is the cotangent:
\[
\Phi(r) = -\frac{GM}{R}\cot\!\left(\frac{r}{R}\right).
\]
That this satisfies the Laplace equation on $S^3$ can be verified directly: in geodesic polar coordinates with metric $ds^2 = dr^2 + R^2 \sin^2(r/R)\,(d\vartheta^2 + \sin^2\!\vartheta\, d\varphi^2)$, the radial Laplacian annihilates $\cot(r/R)$ identically. The Gauss-law normalization is confirmed by
\[
g(r) \cdot 4\pi R^2 \sin^2\!\left(\frac{r}{R}\right) = -\frac{GM}{R^2 \sin^2(r/R)} \cdot 4\pi R^2 \sin^2\!\left(\frac{r}{R}\right) = -4\pi GM.
\]
The flux through geodesic spheres is constant, so $\Phi$ is radially harmonic away from the origin. On the first hemisphere $0<r<\pi R$, the potential vanishes if and only if $r=\pi R/2$: attraction ends at the equator of the three-sphere.

In the limit $r \ll R$, the cotangent expands as
\[
\cot\!\left(\frac{r}{R}\right) = \frac{R}{r} - \frac{r}{3R} - \frac{r^3}{45R^3} - \cdots\,,
\]
yielding
\[
\Phi(r) \approx -\frac{GM}{r} + \frac{GM\,r}{3R^2} + \mathcal{O}\!\left(\frac{r^3}{R^4}\right),
\]
which recovers the standard Newtonian potential as the leading term. The higher-order corrections arise purely from the curvature of $S^3$ and vanish in the decompactification limit $R \to \infty$.

\subsection{Geometric Enhancement of Gravity}

The gravitational acceleration on $S^3$,
\[
g(r) = -\frac{d\Phi}{dr} = -\frac{GM}{R^2 \sin^2(r/R)},
\]
exceeds the flat-space Newtonian acceleration $g_{\text{N}} = -GM/r^2$ by the enhancement factor
\[
\eta(r) \equiv \frac{g_{S^3}(r)}{g_{\text{N}}(r)} = \left[\frac{r/R}{\sin(r/R)}\right]^2.
\]
This factor is precisely the square of the metric distortion of the azimuthal equidistant projection: $\eta\,\mathcal{J}=1$. Gravitational flux through a sphere of geodesic radius $r$ on $S^3$ subtends an area $4\pi R^2 \sin^2(r/R)$, which is smaller than the flat-space area $4\pi r^2$. The flux is concentrated, and gravity is strengthened --- not by invisible matter, but by the curvature of the compact dual spacetime. On $(0,\pi)$ the enhancement is strictly increasing and greater than one. At the potential zero, $\eta(\pi/2)=\pi^2/4$.

Representative windows on the attractive hemisphere: $1.003<\eta<1.004$ at $r/R=0.1$; $1.40<\eta<1.48$ at $r/R=1$. The value $r/R=2$ already lies beyond $\pi/2$ and is not a rotation-curve point.

\subsection{Emergence of Flat Rotation Curves}

For a test particle in circular orbit about a point mass $M$ on the three-sphere chart,
\begin{equation}
\label{eq:s3-circular}
v^2(r)=r\,|g(r)|=\frac{GM}{R}\,f(r/R)=\frac{GM}{r}\,\eta(r/R),
\end{equation}
where $f(x)=x/\sin^2 x$. Flatness of $v$ is constancy of $f$. Outside the luminous disk, where $M(r)\approx M_{\text{tot}}$, the Keplerian factor $GM_{\text{tot}}/r$ declines as $r^{-1}$. The strictly increasing enhancement $\eta$ counteracts this decline. The combined shape is $f$ itself.

On the attractive interval $(0,\pi/2)$ the shape $f$ has a unique critical point, the unique root $x_0$ of $\tan x=2x$ in $(1.14,1.20)$. At that root $f(x_0)=x_0+1/(4x_0)$, and this value $f_0$ is the unique minimum of $f$ on $(0,\pi/2)$:
\[
\frac{135}{100}<f_0<\frac{141}{100}.
\]
That minimum is the geometric origin of a plateau: after the inner Keplerian decline is arrested, $v^2$ sits near $(GM/R)f_0$ throughout a neighbourhood of $r=x_0 R$. One has $x_0<\pi/2<2$, so the attractive patch ends at $r=\pi R/2$, strictly before $r=2R$. The three regions on a point-mass chart are therefore:

\begin{itemize}
  \item \textbf{Inner region} ($x\ll 1$): Newtonian recovery, Keplerian behaviour modified by the rising $M(r)$.
  \item \textbf{Plateau} ($x\sim x_0$): the unique minimum of $f$, with $1.40<f(1)<1.48$ already close to $f_0$.
  \item \textbf{Upturn and reversal} ($x_0<x<\pi/2$): $f$ rises; at $x=\pi/2$ the potential vanishes and attraction ends. The radius $r=2R$ already lies in the repulsive hemisphere.
\end{itemize}

\subsection{The Baryonic Tully-Fisher Relation and the Compactification Radius}

The asymptotic circular velocity at the plateau satisfies
\[
v_{\text{flat}}^2 \sim \frac{GM_{\text{tot}}}{R} \cdot f_0,
\]
where $f_0$ is the unique minimum of $f$ on $(0,\pi/2)$. If the compactification radius scales with the total enclosed mass as
\[
R = \sqrt{\frac{GM_{\text{tot}}}{a_0}},
\]
where $a_0 \approx 1.2 \times 10^{-10}\;\mathrm{m/s^2}$ is an external observational acceleration scale (the baryonic Tully--Fisher / MOND scale), treated here as an input rather than a quantity derived from the dual-rotor algebra, then
\[
v_{\text{flat}}^4 = G M_{\text{tot}} \, a_0 \, f_0^2,
\]
which is the baryonic Tully--Fisher relation as an identity of the chart, given that scaling. The scaling $R \propto M^{1/2}$ implies that more massive galaxies possess larger dual-spacetime patches. A derivation of $a_0$ from the dual-rotor algebra is not claimed.

For the Milky Way ($M_{\text{tot}} \approx 6 \times 10^{10}\,M_\odot$), the same SI stand-ins yield
\[
8\,\mathrm{kpc}<R<9\,\mathrm{kpc},
\]
of which $R\approx 8.3\,\mathrm{kpc}$ is only a reading of the window. The geometric plateau therefore sits at $x_0 R$, in $(9,11)\,\mathrm{kpc}$, inside the optical disk, while the torsional reversal at $\pi R/2$ lies in $(12,15)\,\mathrm{kpc}$. Identifying $a_0$ with a characteristic acceleration on the $S^3$ patch is a proposal.

\subsection{Torsional Reversal and the Finite Range of Gravity}

The cotangent potential crosses zero at $r = \pi R/2$, beyond which the dual-spacetime rotation angle $\theta(r) = r/R$ exceeds $\pi/2$. In the torsional mismatch framework (Section~\ref{sec:gravity}), this crossing triggers a sign change in the torsion scalar $J$: the compact trigonometric contribution from the dual sector overwhelms the non-compact hyperbolic contribution from the usual sector, yielding $J < 0$ and repulsive gravity. The scalar interference factor (Section~\ref{sec:timedilation})
\[
\gamma_s = \cosh\!\left(\frac{\phi}{2}\right)\cos\!\left(\frac{\theta}{2}\right) - \sinh\!\left(\frac{\phi}{2}\right)\sin\!\left(\frac{\theta}{2}\right)
\]
becomes negative when $\theta$ exceeds the critical value satisfying $\tanh(\phi/2)\tan(\theta_c/2) = 1$.

Gravitational attraction is therefore confined to a finite $S^3$ patch of effective radius $\sim \pi R/2$, beyond which a weakly repulsive torsional barrier emerges. This finite range is the compact-space analog of the Newtonian $1/r$ potential extending to infinity: the $S^3$ topology naturally ``walls off'' each gravitational source.

\subsection{The Cosmic Atlas: Galaxy-Scale 3-sphere Patches}

The observable universe is tiled by overlapping $S^3$ patches, each centered on a gravitational concentration (galaxy or cluster), forming a cosmic atlas in the strict differential-geometric sense. Each patch functions as an azimuthal equidistant projection chart: faithful near its center, maximally distorted at its boundary. The large-scale structure of the cosmos is that of a manifold covered by galaxy-scale stickers of azimuthal equidistant maps, glued together at their repulsive boundaries.

If galaxies and clusters are such patches, this patchwork is proposed to account for:
\begin{itemize}
  \item the \textbf{cosmic web}: galaxies are distributed along the overlap regions of neighboring $S^3$ charts, where the attractive tails of adjacent patches constructively interfere;
  \item \textbf{cosmic voids}: the vast underdense regions correspond to the repulsive ($J < 0$) boundary zones between charts, where matter is expelled;
  \item the \textbf{radial acceleration relation}: the tight correlation between observed and baryonic gravitational accelerations follows from the universal $\eta(r)$ enhancement, with residual scatter arising from chart-overlap effects;
  \item the \textbf{absence of dark matter signals} in direct detection experiments --- because dark matter does not exist.
\end{itemize}

\subsection{Numerical Verification and Falsifiable Predictions}

Preliminary particle-based simulations using rotor-ensemble dynamics have successfully reproduced flat rotation curves from baryonic mass distributions alone. Public open-source codes are available at:
\begin{itemize}
  \item \texttt{https://github.com/hypernumbernet/blackhole-simulator}
  \item \texttt{https://github.com/hypernumbernet/dual-spacetime-simulator}
\end{itemize}
Full $N$-body integrations incorporating the exact $S^3$ cotangent potential are in preparation.

Key predictions distinguishing this theory from $\Lambda$CDM and MOND include:
\begin{itemize}
  \item \textbf{Outer rotation curve upturn and reversal}: Isolated galaxies should exhibit a gentle upturn on $(x_0 R,\,\pi R/2)$ and a sign change of gravity at $r=\pi R/2$, not an upturn at $r\gtrsim 2R$. For Milky Way--type masses the plateau lies near $9$--$11\,\mathrm{kpc}$ and the reversal near $12$--$15\,\mathrm{kpc}$, so the signatures are accessible to deep 21-cm HI observations beyond the inner disk rather than exclusively at $\sim 100\,\mathrm{kpc}$. The radius $2R$ is already past reversal.
  \item \textbf{Lensing sign reversal}: Galaxy-galaxy weak lensing should weaken and reverse sign at projected separations exceeding $\pi R/2$, in sharp contrast to the monotonic NFW profiles predicted by $\Lambda$CDM.
  \item \textbf{Cotangent vs.\ MOND interpolation}: The precise form of the radial acceleration relation should follow the cotangent-derived enhancement $g_{\text{obs}} = g_{\text{bar}} \cdot \eta(r)$ rather than the standard MOND interpolation function $\nu(g_{\text{bar}}/a_0)$, with measurable deviations at $g \lesssim a_0/10$.
  \item \textbf{Diversity from $R(M)$}: The observed diversity of rotation curve shapes --- from slowly rising (low-surface-brightness dwarfs with $r_{\text{disk}} \ll R$) to flat (Milky Way-type) to gently declining (high-surface-brightness spirals with $r_{\text{disk}} \sim R$) --- is reproduced by the single relation $R = \sqrt{GM/a_0}$ without additional free parameters.
\end{itemize}

If galaxies are such patches, dark matter is superfluous --- an artifact of projecting the compact $S^3$ geometry of the dual spacetime onto a fictitious flat and infinite background. The cotangent potential, the azimuthal equidistant distortion, and the torsional reversal at the $S^3$ boundary then provide the missing gravity without invisible matter.

\section{Gravitational Engineering: Coherent Excitation of Dual-Spacetime Rotors}
\label{sec:gravitycontrol}

Because $J$ depends only on the relative rotor $\Omega$, an external field that coherently modulates the dual angles $\theta_a$ could in principle reduce $J\to 0$ (shielding or inertial-mass reduction) or drive $J<0$ (repulsion). No electromagnetic coupling, resonance frequency, or laboratory protocol is forced by the dual-rotor algebra. The remainder of this section records a proposed model.

\subsection{Theoretical Coupling of Circularly Polarized Electromagnetic Fields to Dual Rotors}

The electromagnetic field is naturally partitioned across the two spacetimes:

\[
F = F_{\text{usual}} + F_{\text{dual}},
\]

where

\[
F_{\text{usual}} = \mathbf{E} \cdot (iI, iJ, iK) \quad (\hat{p}^2 = +1:\ \text{boost-like, non-compact}),
\]

\[
F_{\text{dual}} = \mathbf{B} \cdot (I, J, K) \quad (\hat{q}^2 = -1:\ \text{rotation-like, compact}).
\]

The electric component $\mathbf{E}$ occupies the usual-spacetime boost generators, while the magnetic component $\mathbf{B}$ occupies the dual-spacetime rotation generators, as in~\eqref{eq:faraday-split}. Laboratory time reversal $(E,B)\mapsto(E,-B)$ is a distinct involution from both duality and the usual--dual parameter swap.

For a circularly polarized wave propagating along $z$ with helicity $\sigma = \pm 1$,

\[
\mathbf{E} = E_0 (\hat{x} \cos(kz - \omega t) + \sigma \hat{y} \sin(kz - \omega t)), \quad \mathbf{B} = \frac{E_0}{c} (-\hat{x} \sin(kz - \omega t) + \sigma \hat{y} \cos(kz - \omega t)).
\]

At every phase that wave is null: $J=0$, $E\cdot B=0$, and $M=E_0^2$ whenever $E_0\neq 0$. Each linear Faraday coefficient is a first harmonic of the phase, so the four-phase means vanish,
\[
\langle F_{\mathrm{usual}}\rangle=\langle F_{\mathrm{dual}}\rangle=0.
\]
There is in particular no mean dual field proportional to $\sigma E_0^2 K$. What is helicity-odd and constant is the bilinear Poynting component $(E\times B)_z=\sigma E_0^2$. Superposition of the wave onto a background torsion does not shift the mean of $J$, and the Faraday mix cannot create $J$ from the wave. A first-order rest sandwich of the dual summand vanishes for every Faraday field, circular waves included. The rest sandwich of the wave itself is the oscillating electric kick, whose four-phase mean vanishes; a velocity along the beam produces a transverse kick with no longitudinal component~\eqref{eq:faraday-lorentz}. Axis projectors $P_{L,R}$ likewise fail to extract helicity.

A helicity-dependent drive of the particle mismatch is therefore not a local scalar of $F$, nor a linear superposition, nor a first-order rest sandwich of $F_{\mathrm{dual}}$, nor an axis-chirality projection, nor a beam-direction Lorentz kick. Any such drive requires extra structure---a preferred axis, a potential, or the particle's dual phase. In the model below, the helicity-odd bilinear $(E\times B)_z=\sigma E_0^2$ is the candidate source of a dual-phase drive, with right-handed polarisation ($\sigma=+1$) reinforcing attractive mismatch ($J>0$) and left-handed polarisation ($\sigma=-1$) opposing it.

The coupling in the working model is written as a path-ordered exponential

\[
R_{\text{dual}} \to \mathcal{P} \exp\!\left( \frac{1}{2} \int (\theta_a \Gamma_a + e \lambda A_a \Gamma_a) \, ds \right),
\]

where the dual generators $\Gamma_a$ mediate the interaction. Near resonance, the dual phase is modelled as

\[
\dot{\theta} = \kappa \sigma E_0^2 \sin(\omega t - \theta + \delta_\sigma),
\]

yielding coherent excitation of the torsional mismatch angle only within this proposed model.

\subsection{Resonance Frequencies and the Role of Superconductors}

The intrinsic dual-spacetime scale is Planckian ($\sim 10^{43}$ Hz). Macroscopic coherence ($N \sim 10^{26}$ particles) redshifts the effective frequency to $10^{12}$--$10^{15}$ Hz. Superconductors further reduce it dramatically via Josephson plasma resonances and collective phase locking, achieving tunable modes in the GHz--THz range accessible to current technology.

\subsection{Proposed Experimental Protocol and Predictions}

A microgram-scale test mass in a high-$Q$ superconducting resonator is illuminated with tunable circularly polarized radiation (0.1--10 THz). Expected signatures include resonant, helicity-dependent changes in apparent weight or inertial mass ($\Delta m/m \gtrsim 10^{-6}$), with left-handed polarization producing reduction and right-handed enhancement. Success at this level scales directly to practical gravitational engineering.

If the model is right, gravity is not merely controllable but chirally controllable, with helicity serving as the switch between attraction and repulsion. That is a conjecture, to be decided in the laboratory.

\section{Baryon Asymmetry: The Geometric Origin from Dual Rotor Time Arrow Fixation}
\label{sec:baryonasymmetry}

We propose that the observed baryon asymmetry $\eta \equiv (n_B - n_{\bar{B}})/n_\gamma \approx 6 \times 10^{-10}$ arises from asymmetric fixation of the time arrow in the dual rotors during the Planck epoch. The construction does not derive $\eta$ from $J$ or $\Omega$; it is a geometric model of the origin of the arrow.

\subsection{The Intrinsic Time Duality and Rotor Parameters}

Recall that each particle encodes a pair of Minkowski spacetimes within the 16-real-dimensional biquaternionic structure: the usual spacetime spanned by $\{j, kI, kJ, kK\}$ with future-directed time basis $j$, and its dual counterpart $\{k, jI, jJ, jK\}$ obtained via right-multiplication by the pseudoscalar $i$, yielding $j i = -k$ and thus a past-directed time basis $k$. This duality operation $X \mapsto X i$ preserves the Minkowski norm (see Sec.~\ref{sec:math}, where the explicit non-commutative computation gives $(X i)^2 = X^2$) but inverts the temporal orientation, imprinting a fundamental chirality: particles (usual-dominant) propagate forward in time, while antiparticles (dual-dominant) would inherently propagate backward if not for dynamical suppression.

The full kinematic evolution is governed by the complete rotor. Torsion-free vacuum in the kinematic sense is $\Omega=1$, i.e.\ $R_{\mathrm{dual}}=R_{\mathrm{usual}}$, not $J=0$. Balanced massive states with $J=0$ and $\Omega\neq 1$ are not particle--antiparticle vacua.

\subsection{Planck-Epoch Fixation of the Time Arrow}

At the Planck scale ($t \sim 10^{-43}$ s, $T \sim 10^{19}$ GeV), quantum gravitational fluctuations---manifest as zero-point oscillations in the dual rotor bivectors---induce a spontaneous symmetry breaking via a phase transition in the vacuum expectation value of the rotor parameters. The biquaternionic basis asymmetry, where $j$ is the sole time-like generator while spatial bases involve cross-terms ($kI$, etc.), selects a unique energetically favorable configuration:

\[
\langle \theta \rangle = \phi + \delta, \quad \delta > 0,
\]
where $\delta$ is the CP-odd chiral phase imprinted at the Planck epoch. We take $\delta$ as a small, free parameter of the framework whose magnitude is to be fixed by observation; the rough estimate $\delta \sim 10^{-5}$ (comparable to the Jarlskog measure of the observed CKM phase) is adopted throughout. This fixation reflects the parity-odd character of the dual map $X \mapsto X i$ itself: the pseudoscalar $i$ flips sign under spatial inversion ($i \to -i$), so any Planck-scale potential that distinguishes $R_\text{usual}$ from $R_\text{dual}$ is necessarily CP-odd in $\delta$ and energetically prefers a definite sign of $\delta$. We do not specify the microscopic form of this potential here; it is governed by quantum gravitational dynamics outside the scope of the present algebraic argument and enters only through the single scalar parameter $\delta > 0$.

To extract the resulting torsional mismatch, we apply the dual conversion $\theta \mapsto \theta i$ established in \S\ref{subsec:dual-rotors}. Because the physical effect of the dual rotor on the dual spacetime is that of a boost (with $(\hat q\, i)^2 = +1$), the over-rotation $\delta$ contributes to $\Omega_\text{biv}$ as a pure boost-direction excess. The pre-existing equilibrium ($\theta = \phi$) yields a bivector whose $i\Gamma_a$ and $\Gamma_a$ coefficients balance on the null cone of the Killing form ($J^{(0)} = 0$, consistent with vacuum torsionlessness); the chiral perturbation breaks this balance asymmetrically. Setting $\delta_a = \delta\, p_a$ along the boost axis $\hat p = \sum_a p_a\, i\Gamma_a$, the linearized mismatch reduces to the boost-only form
\[
\Omega_\text{biv} \approx \sum_a \delta_a\, i\Gamma_a \in \mathfrak{so}(3,1),
\]
structurally identical to the weak-field linearization of \S\ref{sec:gravity}. The Killing-form values established there ($B(i\Gamma_a, i\Gamma_b) = +8\,\delta_{ab}$) then yield directly
\[
J = \frac{1}{2}\sum_a \delta_a^2 = \frac{\delta^2}{2} > 0.
\]
This breaks the particle--antiparticle symmetry at the algebraic core. For usual-dominant modes (matter), the torsional mismatch is positive ($J > 0$) and produces attractive gravitational binding. For dual-dominant modes (antimatter), the usual/dual role assignment is inverted: the inverse conversion $\phi \mapsto -\phi\, i$ maps the chiral phase $\delta$ into the rotation direction rather than the boost direction, so the mismatch becomes $\Omega_\text{biv} \approx \sum_a \delta_a\, \Gamma_a$. The Killing-form value $B(\Gamma_a, \Gamma_b) = -8\,\delta_{ab}$ then gives $J = -\delta^2/2 < 0$, a repulsive torsional layer that dynamically excludes such modes from the cosmological vacuum.

\subsection{Sakharov Conditions and the Order-of-Magnitude Estimate}

The fixation automatically fulfills the three Sakharov conditions through the rotor exponential:

1. Baryon Number Violation: Baryon number $B$ is encoded as the pseudoscalar phase in the dual rotor; specifically, $B \propto \arg(\det R_\text{dual})$. The mismatch $\delta \neq 0$ enables non-conserving processes such as $X \mapsto X i i = -X$ (double time reversal, effectively a baryon flip), so the rates for usual- and dual-dominant modes differ at $O(\delta)$.

2. C and CP Violation: The dual map $i$ is intrinsically chiral (odd under parity, as $i \mapsto -i$ under spatial inversion), and $\delta > 0$ introduces a CP-odd phase $\phi_\text{CP} = \arg(\exp(i \delta \Gamma_a)) \sim \delta$, comparable to the observed CKM phase. This geometric CP violation is universal, permeating all fermion generations without Yukawa-specific tuning.

3. Departure from Thermal Equilibrium: Inflationary expansion drives the system far from equilibrium, allowing the CP-odd phase $\delta$ to be converted into a frozen-in number-density asymmetry as dual-dominant modes annihilate against their usual-dominant partners. A standard freeze-out estimate gives an asymmetry that is quadratic in the chiral phase,
\[
\eta \sim \delta^2,
\]
so that the observed value $\eta \sim 6 \times 10^{-10}$ is reproduced by $\delta \sim 10^{-5}$, in agreement with the CP-odd phase estimate above. The precise numerical coefficient depends on the post-inflationary reheating temperature and is not predicted at this level.

Post-inflation, residual dual-dominant antiparticles annihilate against the baryon surplus, leaving $\eta$ as the frozen relic of the time arrow fixation. Antiparticles do not ``disappear'' but are dynamically excluded: their backward time arrow is incompatible with the forward-directed cosmic expansion fixed by $\delta > 0$.

\subsection{Observational Signatures and Falsifiability}

This mechanism predicts subtle CMB anisotropies from chiral gravitational waves sourced by $\Omega_\text{biv}$, with a tensor-to-scalar ratio $r \sim \delta^2 \sim 10^{-10}$ and a CP-odd parity spectrum detectable by future missions like LiteBIRD. At low energies, it implies a universal lepton asymmetry $\eta_\ell \approx \eta_B$, resolvable by neutrinoless double-beta decay experiments (e.g., LEGEND-200). Non-observation of these signatures would falsify the theory, while confirmation would elevate dual spacetime to the foundational framework for cosmology.

In essence, the baryon asymmetry is not a puzzle but the geometric echo of the universe's choice of time's arrow: forward for matter, forbidden for antimatter. The dual rotors, once fixed, decree that our cosmos is a sanctuary for particles alone.

\section{Discrete Time and Compactified Torsion Scalar}
\label{sec:discretetime}

The rejection of the spacetime continuum hypothesis demands a discrete structure at the fundamental scale. The intrinsic duality encoded in the biquaternion algebra—particularly the pseudoscalar $i$ and the dual map $X \mapsto Xi$ that reverses the temporal arrow while preserving the Minkowski norm—compactifies the dual spacetime intrinsically, while leaving the usual spacetime non-compact.

We therefore discretize only the dual-spacetime rotation angle:
\[
\theta \in \frac{2\pi}{N} \mathbb{Z},
\]
where $N \gg 1$ is a very large integer determining the compactification scale of the dual spacetime (typically proportional to the particle’s rest mass in Planck units). The usual-spacetime rapidity remains continuous:
\[
\phi \in \mathbb{R}.
\]

This asymmetry is required for theoretical consistency:
\begin{itemize}
  \item The dual spacetime is elliptic ($\hat{q}^2 = -1$), inherently compact, and its discretization directly reflects the $4\pi$-periodic topology of $S^3$ parameterized by $R_{\text{dual}}$. This provides the finite phase volume essential for resolving the Newton–Wigner localization problem and yielding natural quantum discreteness.
  \item The usual spacetime is hyperbolic ($\hat{p}^2 = +1$), intrinsically non-compact, so continuity of $\phi$ preserves the full Lorentz group in the macroscopic limit.
\end{itemize}

The torsional mismatch rotor $\Omega = R_{\text{usual}}^\dagger R_{\text{dual}}$ thus combines a continuous usual component with a discrete dual component. The scalar $J$ is bounded by $|J|\le\Jbound$ only on admissible configurations. Discretizing the dual angle alone does not imply the bound: a principal-branch condition without non-negative rapidities admits counterexamples, and the continuous usual-sector rapidity is unbounded. Continuum-induced singularities are therefore excluded only after the admissible cone is imposed. On the discrete torus $(\mathbb{Z}/N\mathbb{Z})^6$ the rotor image is finite. Nonzero winding of a scaled lattice configuration is incompatible with real-scale admissibility. When $4\mid N$ the lattice attains the extremals $\Jnorm=\pm 1$; otherwise the admissible bound is strict.

A critical issue arises in the effective Lorentz factor across torsional layers, which typically appears as a synchronization function of the form
\[
\gamma_{\text{eff}} \propto D_N(\delta\theta) \equiv \frac{\sin(N \delta\theta / 2)}{N \sin(\delta\theta / 2)},
\]
where $\delta\theta$ is the torsional mismatch angle and $D_N$ is the (normalized) Dirichlet kernel. The denominator vanishes precisely on the lattice $\delta\theta \in 2\pi \mathbb{Z}$.

At every such point the numerator $\sin(N \cdot k\pi)$ also vanishes (for any integer $N$, prime or composite), so the apparent singularity is in fact a removable $0/0$ of Dirichlet-kernel type. The standard L'H\^{o}pital evaluation gives the finite limit
\[
\lim_{\delta\theta \to 2\pi k} D_N(\delta\theta) = \pm 1 \qquad (k \in \mathbb{Z}),
\]
so the synchronization function is bounded on the entire compact lattice and the classical continuum limit $N \to \infty$ is recovered without pathology.

Thus the Dirichlet kernel is bounded for any large positive integer $N$; primality is not required. Filling the removable singularities by the L'H\^{o}pital value yields $|D_N(\theta)|\le 1$, with $D_N=\pm 1$ at every lattice point. The relation between $N$ and particle mass spectra, and the zeta-regularized product $\prod_p p\to 4\pi^2$, remain open. The geometric reading of $4\pi^2$ as the square of the dual-rotor period $2\pi$ is a consistency remark, not a derivation of mass spectra.

\section{Dual-Rotor Dynamics and the Geometric Origin of the de Broglie Phase}
\label{sec:dualrotordynamics}

The mismatch scalar $J$ is proposed as a gravitational action density in the classical field theory. It also admits a particle-intrinsic reading as a potential for the internal angles of individual massive particles. That reading, and the de Broglie identification below, are proposals.

For small mismatch angles, the Killing form invariant expands as
\[
J \approx \frac{1}{2} \sum_{a=1}^3 \delta\phi_a^2,
\]
where $\delta\phi_a = \phi_a - \theta_a$ is the torsional mismatch angle in each boost-rotation plane. Incorporating the rest mass $m$ as a rigidity constant of the dual sector is a modelling choice, not a derivation of $m$ from the algebra:
\[
J = \frac{1}{2} m \sum_{a=1}^3 (\phi_a - \theta_a)^2.
\]

We consider the written effective action for the intrinsic rotor angles of a single particle:
\[
S_{\text{particle}} = \int \left[ \frac{1}{2} \sum_{a=1}^3 \dot{\phi}_a^2 - \frac{1}{2} \sum_{a=1}^3 \dot{\theta}_a^2 - \frac{m}{2} \sum_{a=1}^3 (\phi_a - \theta_a)^2 \right] dt.
\]
The opposite signs of the kinetic terms are part of the written model; they are not deduced from unitarity of the rotors.

The Euler--Lagrange equations of this Lagrangian are
\[
\ddot\phi=-m(\phi-\theta),\qquad \ddot\theta=-m(\phi-\theta).
\]
The mismatch $\delta=\phi-\theta$ is therefore free: $\ddot\delta=0$. The written action does not yield the oscillator $\ddot\delta+m\delta=0$. A neighbouring system sometimes written alongside it,
\[
\ddot\phi=m(\phi-\theta),\qquad -\ddot\theta=m(\phi-\theta),
\]
makes the mismatch runaway, $\ddot\delta=2m\delta$, and is likewise not an oscillator.

An oscillator does appear if both kinetic signs are taken positive. The modified density $L=\tfrac12\dot\phi^2+\tfrac12\dot\theta^2-\tfrac m2(\phi-\theta)^2$ has Euler--Lagrange equations $\ddot\phi=-m(\phi-\theta)$ and $\ddot\theta=m(\phi-\theta)$, hence
\[
\ddot\delta+2m\delta=0.
\]
This is a separate working model, not the written action. Its frequency is $\sqrt{2m}$ (in natural units $\hbar=c=1$). For a free particle in uniform motion the usual rapidity is modelled as evolving at constant rate $\dot{\phi}_a=\gamma v_a/c$. The general solution is
\[
\delta\phi_a(t) = A_a \sin(\sqrt{2m}\, t + \psi_a).
\]
In the low-energy, non-relativistic regime, rapid Compton-frequency oscillations are assumed to average to zero, leaving a slow linear accumulation of phase proportional to the proper time. Restoring $\hbar$ and $c$, the accumulated phase is conjectured to be
\[
\int \delta\phi_a \, dt \propto \frac{E t - \mathbf{p} \cdot \mathbf{x}}{\hbar},
\]
the de Broglie phase of the matter wave. Collective averaging of this particle-intrinsic action over many particles is proposed to recover a macroscopic integral of $J$. That averaging is not an identification of $J$ with the Weitzenb\"ock density $T$. The written action does not furnish a harmonic bridge to quantum theory; the same-sign model is the only oscillator among the three systems compared here.

\section{Towards Quantum Gravity and Full Unification}
\label{sec:quantumgravity}

The classical formulation presented thus far treats gravity as torsional mismatch between the particle-intrinsic usual and dual spacetimes, encoded in $\mathbb{H}\otimes_{\mathbb{R}}\mathbb{H}\cong\mathrm{Cl}(3,1)$. The proposed action uses the parameter-space scalar $J$; it is not identified with the Einstein--Hilbert integral.

The dual rotor $R_\text{dual}=\exp(\theta/2\,\hat q)$ is the spin-$1/2$ rotation operator of $\mathrm{SU}(2)$. Dual-sector phases $\theta$ therefore parameterize quantum states. Represented on $\mathbb{C}^2$ by $\exp\bigl(\sum_a(\theta_a/2)(-i\sigma_a)\bigr)$, the dual rotor is unitary with determinant $1$, hence an element of $\mathrm{SU}(2)$. Due to the compactness of the dual spacetime ($\theta\bmod 4\pi$), these phases live on a compact manifold: each particle's state $|\psi\rangle$ corresponds to the unitary phase part of $R_\text{dual}$. A Hilbert-space completion of a $\mathrm{Cl}(3,1)$ minimal left ideal is not constructed here.

Position and momentum operators, and uncertainty relations among them, are not derived. Singularities are excluded on the admissible cone by the bound of Section~\ref{sec:gravity}; the quasi-horizon of Section~\ref{subsec:quasi-horizon} is a reading of that bound on the saturated Schwarzschild chart. An explicit Hilbert space of $\mathrm{Cl}(3,1)$ spinors, and a complete quantization of the mismatch dynamics, remain open.

Quantum gravity is therefore not a separate theory: it is the dual-rotor kinematics of the same 16-dimensional algebra, still incomplete as a Hilbert-space construction.

\section{Conclusions}
\label{sec:conclusions}

The century that began with Einstein taught us to speak of gravity as the curvature of a continuum. That speech is exact as a macroscopic bookkeeping, and it remains so. What it does not say is why a continuum should curve, nor why the same phenomenon should appear as inertia under acceleration. The present work has refused the continuum as a primitive. There is no autonomous stage. Each massive particle carries, in the sixteen-real-dimensional biquaternion algebra isomorphic to $\mathrm{Cl}(3,1)$, a pair of Minkowski sectors---the usual and the dual. A map between them reverses the arrow of time and leaves the Minkowski norm untouched. Gravity is the torsion of that pair: the mismatch of two rotors that belong to the particle itself.

The relative rotor is the identity precisely when the two rotors coincide. A Lorentz-invariant scalar $J$ measures the mismatch; it is bounded on the admissible cone, and vanishing $J$ does not force the rotors to agree. On every static radial-boost chart, $J$ stands in a closed relation with the Weitzenb\"ock torsion $T$ of teleparallel gravity, of which the exterior Schwarzschild geometry is the unique vacuum case. The two agree only in the flat limit. A real radial boost cannot produce repulsion; repulsive layers belong to the elliptic sector and are of bounded depth. Where the two rotors share a common scale, $J$ vanishes while their interference still alternates, so the layers of strong-field gravity are ordered by that interference and not by $J$. The same algebra houses electromagnetism: Faraday coefficients occupy the six-generator space of the dual rotor, and the sandwich increment is a Lorentz 4-force. On a three-sphere chart the circular speed is $v^2=(GM/R)f(r/R)$, the shape $f$ has a unique minimum on $(0,\pi/2)$, and attraction ends at $\pi R/2$, before $r=2R$. None of this requires a background manifold.

We have proposed that the layered phase is the strong force; that galaxies are three-sphere patches of the dual sector, and that this compactness therefore accounts for flat rotation curves; that coherent excitation of the dual sector might engineer gravity; that early dual dynamics sources baryon asymmetry; that Maxwell's equations, a Weinberg angle, and a Higgs reading of rotor rigidity can be recovered from the same six-space; that a mismatch oscillator is the de Broglie phase; and that a Hilbert-space quantization of the dual rotor completes the theory. These are readings of the geometry. Controllability of gravity is a conjecture.

What can already be asked of the sky is modest and sharp. In place of a true horizon one expects a finite redshift floor. At the centre of a spherically symmetric torsion star, gravitational time dilation should vanish. If a galaxy is a three-sphere patch, its rotation curve plateaus near $r\sim x_0 R$ with $x_0\in(1.14,1.20)$ and reverses at $\pi R/2$, rather than by unseen mass; for Milky-Way stand-in mass the plateau sits near $9$--$11\,\mathrm{kpc}$ and the reversal near $12$--$15\,\mathrm{kpc}$. The ratio of successive barrier radii is a pure number, the same at nuclear and at stellar scales, because the characteristic length cancels from it.

If this ontology is sound, then the geometry of the world is not prior to matter. It is an attribute of particles, as spin and charge are attributes, and the apparent curvature of the heavens is the collective torsion of their dual rotors. The continuum remains a useful fiction of the large. What remains is not to curve an empty stage more ingeniously, but to learn, from the algebra already in hand, how the paired rotors of matter compose the world we see.

\clearpage
\appendix
\section{Multiplication Table}
\label{app:mult-table}

\begin{table}[ht]
  \begin{center}
    \caption{Biquaternion Multiplication Table}
    \small
    \begin{tabular}{r|rrrrrrrrrrrrrrr}
      &$j$&$kI$&$kJ$&$kK$&$iI$&$iJ$&$iK$&$I$&$J$&$K$&$k$&$jI$&$jJ$&$jK$&$i$ \\ \hline
      j&$-1$&$+iI$&$+iJ$&$+iK$&$-kI$&$-kJ$&$-kK$&$+jI$&$+jJ$&$+jK$&$+i$&$-I$&$-J$&$-K$&$-k$ \\ \hline
      kI&$-iI$&$+1$&$-K$&$+J$&$-j$&$+jK$&$-jJ$&$-k$&$+kK$&$-kJ$&$-I$&$+i$&$-iK$&$+iJ$&$+jI$ \\ \hline
      kJ&$-iJ$&$+K$&$+1$&$-I$&$-jK$&$-j$&$+jI$&$-kK$&$-k$&$+kI$&$-J$&$+iK$&$+i$&$-iI$&$+jJ$ \\ \hline
      kK&$-iK$&$-J$&$+I$&$+1$&$+jJ$&$-jI$&$-j$&$+kJ$&$-kI$&$-k$&$-K$&$-iJ$&$+iI$&$+i$&$+jK$ \\ \hline
      iI&$+kI$&$+j$&$-jK$&$+jJ$&$+1$&$-K$&$+J$&$-i$&$+iK$&$-iJ$&$-jI$&$-k$&$+kK$&$-kJ$&$-I$ \\ \hline
      iJ&$+kJ$&$+jK$&$+j$&$-jI$&$+K$&$+1$&$-I$&$-iK$&$-i$&$+iI$&$-jJ$&$-kK$&$-k$&$+kI$&$-J$ \\ \hline
      iK&$+kK$&$-jJ$&$+jI$&$+j$&$-J$&$+I$&$+1$&$+iJ$&$-iI$&$-i$&$-jK$&$+kJ$&$-kI$&$-k$&$-K$ \\ \hline
      I&$+jI$&$-k$&$+kK$&$-kJ$&$-i$&$+iK$&$-iJ$&$-1$&$+K$&$-J$&$+kI$&$-j$&$+jK$&$-jJ$&$+iI$ \\ \hline
      J&$+jJ$&$-kK$&$-k$&$+kI$&$-iK$&$-i$&$+iI$&$-K$&$-1$&$+I$&$+kJ$&$-jK$&$-j$&$+jI$&$+iJ$ \\ \hline
      K&$+jK$&$+kJ$&$-kI$&$-k$&$+iJ$&$-iI$&$-i$&$+J$&$-I$&$-1$&$+kK$&$+jJ$&$-jI$&$-j$&$+iK$ \\ \hline
      k&$-i$&$-I$&$-J$&$-K$&$+jI$&$+jJ$&$+jK$&$+kI$&$+kJ$&$+kK$&$-1$&$-iI$&$-iJ$&$-iK$&$+j$ \\ \hline
      jI&$-I$&$-i$&$+iK$&$-iJ$&$+k$&$-kK$&$+kJ$&$-j$&$+jK$&$-jJ$&$+iI$&$+1$&$-K$&$+J$&$-kI$ \\ \hline
      jJ&$-J$&$-iK$&$-i$&$+iI$&$+kK$&$+k$&$-kI$&$-jK$&$-j$&$+jI$&$+iJ$&$+K$&$+1$&$-I$&$-kJ$ \\ \hline
      jK&$-K$&$+iJ$&$-iI$&$-i$&$-kJ$&$+kI$&$+k$&$+jJ$&$-jI$&$-j$&$+iK$&$-J$&$+I$&$+1$&$-kK$ \\ \hline
      i&$+k$&$-jI$&$-jJ$&$-jK$&$-I$&$-J$&$-K$&$+iI$&$+iJ$&$+iK$&$-j$&$+kI$&$+kJ$&$+kK$&$-1$ \\ \hline
    \end{tabular}
  \end{center}
\end{table}
\begin{table}[ht]
  \begin{center}
    \caption{Cl(3,1) Variant Multiplication Table ($e_0, e_1, e_2, e_3 \mapsto 0,1,2,3$)}
    \scriptsize
    \begin{tabular}{r|rrrrrrrrrrrrrrr}
      &$0$&$1$&$2$&$3$&$01$&$02$&$03$&$32$&$13$&$21$&$123$&$032$&$013$&$021$&$0123$\\\hline
      $0$&$-$&$+01$&$+02$&$+03$&$-1$&$-2$&$-3$&$+032$&$+013$&$+021$&$+0123$&$-32$&$-13$&$-21$&$-123$\\
      $1$&$-01$&$+$&$-21$&$+13$&$-0$&$+021$&$-013$&$-123$&$+3$&$-2$&$-32$&$+0123$&$-03$&$+02$&$+032$\\
      $2$&$-02$&$+21$&$+$&$-32$&$-021$&$-0$&$+032$&$-3$&$-123$&$+1$&$-13$&$+03$&$+0123$&$-01$&$+013$\\
      $3$&$-03$&$-13$&$+32$&$+$&$+013$&$-032$&$-0$&$+2$&$-1$&$-123$&$-21$&$-02$&$+01$&$+0123$&$+021$\\
      $01$&$+1$&$+0$&$-021$&$+013$&$+$&$-21$&$+13$&$-0123$&$+03$&$-02$&$-032$&$-123$&$+3$&$-2$&$-32$\\
      $02$&$+2$&$+021$&$+0$&$-032$&$+21$&$+$&$-32$&$-03$&$-0123$&$+01$&$-013$&$-3$&$-123$&$+1$&$-13$\\
      $03$&$+3$&$-013$&$+032$&$+0$&$-13$&$+32$&$+$&$+02$&$-01$&$-0123$&$-021$&$+2$&$-1$&$-123$&$-21$\\
      $32$&$+032$&$-123$&$+3$&$-2$&$-0123$&$+03$&$-02$&$-$&$+21$&$-13$&$+1$&$-0$&$+021$&$-013$&$+01$\\
      $13$&$+013$&$-3$&$-123$&$+1$&$-03$&$-0123$&$+01$&$-21$&$-$&$+32$&$+2$&$-021$&$-0$&$+032$&$+02$\\
      $21$&$+021$&$+2$&$-1$&$-123$&$+02$&$-01$&$-0123$&$+13$&$-32$&$-$&$+3$&$+013$&$-032$&$-0$&$+03$\\
      $123$&$-0123$&$-32$&$-13$&$-21$&$+032$&$+013$&$+021$&$+1$&$+2$&$+3$&$-$&$-01$&$-02$&$-03$&$+0$\\
      $032$&$-32$&$-0123$&$+03$&$-02$&$+123$&$-3$&$+2$&$-0$&$+021$&$-013$&$+01$&$+$&$-21$&$+13$&$-1$\\
      $013$&$-13$&$-03$&$-0123$&$+01$&$+3$&$+123$&$-1$&$-021$&$-0$&$+032$&$+02$&$+21$&$+$&$-32$&$-2$\\
      $021$&$-21$&$+02$&$-01$&$-0123$&$-2$&$+1$&$+123$&$+013$&$-032$&$-0$&$+03$&$-13$&$+32$&$+$&$-3$\\
      $0123$&$+123$&$-032$&$-013$&$-021$&$-32$&$-13$&$-21$&$+01$&$+02$&$+03$&$-0$&$+1$&$+2$&$+3$&$-$\\
    \end{tabular}
  \end{center}
\end{table}
These are indeed the same results and are isomorphic.

\end{document}
